% Options for packages loaded elsewhere
\PassOptionsToPackage{unicode}{hyperref}
\PassOptionsToPackage{hyphens}{url}
%
\documentclass[
]{article}
\usepackage{amsmath,amssymb}
\usepackage{lmodern}
\usepackage{iftex}
\ifPDFTeX
  \usepackage[T1]{fontenc}
  \usepackage[utf8]{inputenc}
  \usepackage{textcomp} % provide euro and other symbols
\else % if luatex or xetex
  \usepackage{unicode-math}
  \defaultfontfeatures{Scale=MatchLowercase}
  \defaultfontfeatures[\rmfamily]{Ligatures=TeX,Scale=1}
\fi
% Use upquote if available, for straight quotes in verbatim environments
\IfFileExists{upquote.sty}{\usepackage{upquote}}{}
\IfFileExists{microtype.sty}{% use microtype if available
  \usepackage[]{microtype}
  \UseMicrotypeSet[protrusion]{basicmath} % disable protrusion for tt fonts
}{}
\makeatletter
\@ifundefined{KOMAClassName}{% if non-KOMA class
  \IfFileExists{parskip.sty}{%
    \usepackage{parskip}
  }{% else
    \setlength{\parindent}{0pt}
    \setlength{\parskip}{6pt plus 2pt minus 1pt}}
}{% if KOMA class
  \KOMAoptions{parskip=half}}
\makeatother
\usepackage{xcolor}
\setlength{\emergencystretch}{3em} % prevent overfull lines
\providecommand{\tightlist}{%
  \setlength{\itemsep}{0pt}\setlength{\parskip}{0pt}}
\setcounter{secnumdepth}{-\maxdimen} % remove section numbering
\ifLuaTeX
  \usepackage{selnolig}  % disable illegal ligatures
\fi
\IfFileExists{bookmark.sty}{\usepackage{bookmark}}{\usepackage{hyperref}}
\IfFileExists{xurl.sty}{\usepackage{xurl}}{} % add URL line breaks if available
\urlstyle{same} % disable monospaced font for URLs
\hypersetup{
  hidelinks,
  pdfcreator={LaTeX via pandoc}}

\author{Dietmar Gerald SCHRAUSSER}
\date{12.04.2026}

\begin{document}

\hypertarget{console-applications-for-mathematical-functions-eeg-signal-analysis-string-manipulation-and-utilities.}{%
\section{Console Applications for Mathematical Functions, EEG Signal
Analysis, String Manipulation and
Utilities.}\label{console-applications-for-mathematical-functions-eeg-signal-analysis-string-manipulation-and-utilities.}}

\textbf{Dietmar G. Schrausser}\\
\href{https://orcid.org/0000-0002-4924-8280}{orcid.org/0000-0002-4924-8280}

Karl-Franzens University, Graz, Austria

\hypertarget{overview}{%
\subsection{Overview}\label{overview}}

Applications for calculating (1) \emph{distribution} functions and
\emph{bootstrap} estimators (s. Tab. 1; c.f. Schrausser,
\href{https://doi.org/10.5281/zenodo.7664141}{2023a},
\href{https://doi.org/10.31234/osf.io/rvzxa}{2024a}), (2) various
parameters within the framework of psychophysiological \emph{EEG}
measurements (c.f. Schrausser,
\href{https://doi.org/10.5281/zenodo.10701349}{2024b}), (3)
\emph{integral} and \emph{interpolation} (Schrausser,
\href{https://doi.org/10.5281/zenodo.7655056}{2023b}) and (4)
\emph{matrix} parameters (Schrausser,
\href{https://doi.org/10.5281/zenodo.7655046}{2023c}), s. Tab. 2.
Further applications are (5) \emph{string-} and \emph{number-}
\emph{transformation tools} (Schrausser,
\href{https://doi.org/10.5281/zenodo.7653790}{2023d}) and finally (6)
various \emph{console tools} (Schrausser,
\href{https://doi.org/10.5281/zenodo.7655239}{2023e}), see Tab. 3. For a
full \emph{BibTeX} source of the reference list, see Schrausser
(\href{https://doi.org/10.5281/zenodo.19513238}{2026}).

The applications are written in \emph{ANSI-C} (c.f. Gerlach,
\href{https://doi.org/10.1007/978-3-662-59246-5_4}{2019}; Joyce,
\href{https://doi.org/10.1007/978-1-4842-5064-8}{2019}; Gonzalez-Morris
\& Horton, \href{https://doi.org/10.1007/979-8-8688-0149-5}{2024}),
executables were compiled for \emph{MS-DOS} (c.f. Kaier,
\href{https://doi.org/10.1007/978-3-322-89035-1_1}{1990}; Herrmann,
\href{https://doi.org/10.1007/978-3-322-94365-1_16}{2001}).

Certain applications are implemented in the Windows Tools (i)
\texttt{FunktionWin} (see Schrausser,
\href{https://doi.org/10.5281/zenodo.7651661}{2023f},
\href{https://doi.org/10.5281/zenodo.17880113}{2025c}), (ii)
\texttt{ThetaWin} (see Schrausser,
\href{https://www.academia.edu/81800920}{2009},
\href{https://doi.org/10.5281/zenodo.7659264}{2023g},
\href{https://doi.org/10.5281/zenodo.17880113}{2025a}) and (iii)
\texttt{ATINA} (s. Schrausser,
\href{https://doi.org/10.5281/zenodo.19422127}{2006}). For better
comparability, the \emph{original} German-language terms are presented
comparatively in \emph{italic} text.

\hypertarget{consoleapp_distributionfunctions}{%
\subsection{1.
ConsoleApp\_DistributionFunctions}\label{consoleapp_distributionfunctions}}

Console applications for \emph{distribution} functions (c.f. Schrausser,
\href{https://doi.org/10.5281/zenodo.7664141}{2023a},
\href{https://doi.org/10.31234/osf.io/rvzxa}{2024a}) implemented in
\texttt{FunktionWin} (see Schrausser,
\href{https://doi.org/10.5281/zenodo.7651661}{2023f},
\href{https://doi.org/10.5281/zenodo.17880113}{2025c}).

The following functions are available:

(1.1.-1.2.) \emph{Binomial} distribution (de Moivre,
\href{https://doi.org/10.1098/rstl.1710.0018}{1711}; Bernoulli,
\href{https://www.e-rara.ch/zut/doi/10.3931/e-rara-9001}{1713}) as a
discrete distribution of probability values from frequency ratios of two
groups, c.f. Collani and Dräger
(\href{https://doi.org/10.1007/978-1-4612-0215-8}{2001}), Philippou and
Antzoulakos (\href{https://doi.org/10.1007/978-3-662-69359-9_71}{2025}),
further e.g.~Mulholland and Jones
(\href{https://doi.org/10.1007/978-1-4899-6507-3_3}{1968a}), Altham
(\href{http://www.jstor.org/stable/2346943}{1978}), Goel and Rodriguez
(\href{http://www.jstor.org/stable/2689344}{1987}), Vellaisamy and
Punnen (\href{http://www.jstor.org/stable/3215739}{2001}), García-García
et al. (\href{https://doi.org/10.3390/math10152680}{2022}).

(1.3.) Effect size, \emph{Cohen's} \(d\) (Cohen,
\href{https://doi.org/10.1016/C2013-0-10517-X}{1977},
\href{https://doi.org/10.4324/9780203771587}{1988}, p.20, p.49,
\href{https://doi.org/10.1037/0033-2909.112.1.155}{1992}), an effect
size measure that represents the magnitude of the difference between two
group means in standard deviation units. Further works on this topic are
given by e.g.~Wilcox
(\href{http://www.jstor.org/stable/20157436}{2006}), Miller et
al.~(\href{https://doi.org/10.1016/j.jval.2010.10.013}{2011}), Fritz et
al.~(\href{https://doi.org/10.1037/a0024338}{2012}), Janczyk and Pfister
(\href{https://doi.org/10.1007/978-3-642-34825-9_7}{2013}), Peng and
Chen (\href{https://www.jstor.org/stable/26594399}{2014}), Li
(\href{https://doi.org/10.3758/s13428-015-0667-z}{2016}), Schäfer and
Schwarz (\href{https://doi.org/10.3389/fpsyg.2019.00813}{2019}), Bowring
et al. (\href{https://doi.org/10.1016/j.neuroimage.2020.117477}{2021}),
Groß and Möller
(\href{https://doi.org/10.1007/s42519-023-00323-w}{2023},
\href{https://doi.org/10.1007/s00362-023-01527-9}{2024}) and Brandmaier
(\href{https://doi.org/10.31234/osf.io/n9ta7_v1}{2025}).

(1.4.) Fisher-Snedecor \(F\)-distribution (Fisher,
\href{https://repository.rothamsted.ac.uk/item/8w2q9/on-a-distribution-yielding-the-error-functions-of-several-well-known-statistics}{1924}),
a continuous probability distribution for determining significance
levels in multigroup and factorial designs, c.f. also David
(\href{http://www.jstor.org/stable/2332676}{1949}), Patnaik
(\href{http://www.jstor.org/stable/2332542}{1949}), Zinger
(\href{http://www.jstor.org/stable/2985223}{1964}), Cacoullos
(\href{http://www.jstor.org/stable/2282687}{1965}), Stange
(\href{https://doi.org/10.1007/978-3-642-85602-0_11}{1970}), Saunders
and Moran (\href{https://doi.org/10.2307/3213414}{1978}), Herrmann
(\href{https://doi.org/10.1007/978-3-322-96320-8_21}{1984}), Selvin
(\href{https://doi.org/10.1002/9781118445112.stat05856}{2014}), Brereton
(\href{https://doi.org/10.1002/cem.2734}{2015}) or Chattamvelli and
Shanmugam (\href{https://doi.org/10.1007/978-3-031-02435-1_7}{2021}).

(1.5.) \emph{Fisher's exact test} (Fisher,
\href{https://doi.org/10.2307/2340521}{1922}; s. Altham,
\href{http://www.jstor.org/stable/2984209}{1969}), \emph{exact}
\emph{hypergeometric} 4-field test, a nonparametric statistical
procedure for significance testing of frequency distributions in 2×2
designs, see in this context e.g.~Upton
(\href{http://www.jstor.org/stable/2982890}{1992}), Camilli
(\href{https://doi.org/10.1007/BF02301418}{1995}), Hershberger
(\href{https://doi.org/10.1002/9781118445112.stat06165}{2014}), Sprent
(\href{https://doi.org/10.1007/978-3-662-69359-9_693}{2025}).

(1.6.) \emph{Fisher} \(Z\)-transformation (Fisher,
\href{https://doi.org/10.2307/2331838}{1915}), to adjust the
distribution of \emph{product moment} correlation coefficients \(r\) to
a \emph{normal} distribution in order to perform valid significance
tests, interesting discussions on the topic are given by Hjelm and
Norris (\href{http://www.jstor.org/stable/20156574}{1962}), Mendoza
(\href{https://doi.org/10.1007/BF02294830}{1993}), Bond and Richardson
(\href{https://doi.org/10.1007/BF02295945}{2004}), Carbonell et al.
(\href{https://doi.org/10.1016/j.spl.2008.11.007}{2009}) or Yang et al.
(\href{https://doi.org/10.1007/s12209-013-1978-8}{2013}).

(1.7.-1.8.) \emph{Gamma} function \(\Gamma\) (Bernoulli,
\href{https://commons.m.wikimedia.org/wiki/File:DanielBernoulliLetterToGoldbach-1729-10-06.jpg}{1729})
to extend the factorial to noninteger arguments (c.f. Borwein \&
Corless, \href{https://www.jstor.org/stable/48663320}{2018}). The
\emph{gamma} integral, or \emph{gamma} value, is also a \emph{central}
parameter of the \(F\), \(t\), and \(\chi ²\) distributions implemented
here, s. e.g. Gronwall
(\href{http://www.jstor.org/stable/1967180}{1918}), Rasch
(\href{http://www.jstor.org/stable/1968254}{1931}), Davis
(\href{http://www.jstor.org/stable/2309786}{1959}), Magnus et al.
(\href{https://doi.org/10.1007/978-3-662-11761-3_1}{1966a}), Saunders
and Moran (\href{https://doi.org/10.2307/3213414}{1978}), Batir
(\href{https://doi.org/10.1016/j.exmath.2007.10.001}{2008}), Koepf
(\href{https://doi.org/10.1007/978-1-4471-6464-7_1}{2014}).

(1.9.) \emph{Hypergeometric} distribution (Wallis,
\href{https://books.google.com/books?id=Z5w_AAAAcAAJ}{1656}) for the
calculation of significance levels for frequencies in 2×2 designs, see
\emph{Fisher's} \emph{exact} \emph{test} above, c.f. Berggren et al.
(\href{https://doi.org/10.1007/978-1-4757-3240-5_11}{2000}) or Stedall
(\href{https://doi.org/10.1016/B978-044450871-3/50083-8}{2005}), s.
further e.g.~Magnus et al.
(\href{https://doi.org/10.1007/978-3-662-11761-3_2}{1966b}), Dutka
(\href{http://www.jstor.org/stable/41133728}{1984}) or Shuster
(\href{https://doi.org/10.1002/9781118445112.stat04869}{2014})

(1.10.) \emph{Geometric} distribution (de Montmort,
\href{https://books.google.com/books?id=PII_AAAAcAAJ}{1713}) calculates
the probability that event \(A\) with probability \(p_{A}\) occurs at
least once in \(r + 1\) attempts, c.f. Sreehari
(\href{http://www.jstor.org/stable/3213738}{1983}), Zijlstra
(\href{https://doi.org/10.2307/3213595}{1983}) or Thomopoulos
(\href{https://doi.org/10.1007/978-3-319-65112-5_15}{2017}).

(1.11.) \emph{Poisson} distribution (Poisson,
\href{https://gallica.bnf.fr/ark:/12148/bpt6k110193z/f218.image}{1837})
is generally right-skewed, but approaches with (i) increasing expected
number of events \(\lambda\) the \emph{normal} distribution (c.f. Chung,
\href{https://doi.org/10.1007/978-1-4757-3973-2_7}{1974}), with (ii)
large \(n\) and small \(p\) the \emph{binomial} distribution, a more
detailed insight into this method provide e.g.~Rao and Chakravarti
(\href{http://www.jstor.org/stable/3001466}{1956}), Crow
(\href{http://www.jstor.org/stable/2333201}{1958}), Haight
(\href{https://books.google.com/books?id=l8Y-AAAAIAAJ}{1967}), Simonton
(\href{http://www.jstor.org/stable/284821}{1978}), Jolicoeur
(\href{https://doi.org/10.1007/978-1-4615-4777-8_19}{1999}),
Finkelshtein et
al.~(\href{https://doi.org/10.1007/s11009-025-10171-9}{2025}) or Zhang
(\href{https://doi.org/10.54254/2754-1169/2025.BJ24761}{2025}).

(1.12.) Student's \(t\)-distribution (Lüroth,
\href{https://doi.org/10.1002/asna.18760871402}{1876}; Gosset,
\href{https://doi.org/10.2307/2331554}{1908}) to approximate
significance levels if sample sizes \(n\) are small, it converges to the
\emph{normal} distribution as \(n\) increases (s. Wang,
\href{https://doi.org/10.1016/j.spl.2011.03.037}{2011}), c.f. further
Cacoullos (\href{http://www.jstor.org/stable/2282687}{1965}), Yang et
al. (\href{https://doi.org/10.1016/j.jmva.2006.11.003}{2007}), Fonseca
et al. (\href{http://www.jstor.org/stable/20441467}{2008}), Peña et al.
(\href{https://doi.org/10.1007/978-3-540-85636-8_15}{2009}), Grigelionis
(\href{https://doi.org/10.1007/978-3-642-31146-8}{2012}), Brereton
(\href{https://doi.org/10.1002/cem.2734}{2015}), Li and Nadarajah
(\href{https://doi.org/10.1007/s00181-018-1570-0}{2020}), Reyes et al.
(\href{https://doi.org/10.3390/sym13122444}{2021}), Kirkby et al.
(\href{https://doi.org/10.48550/arXiv.1912.01607}{2024}) and Edelman
(\href{https://doi.org/10.48550/arXiv.2508.13226}{2025}).

(1.13.) \(\chi ²\)-distribution (Helmert,
\href{https://gdz.sub.uni-goettingen.de/id/PPN599415665_0021}{1876}) can
be derived from the \emph{normal} distribution (c.f. Behnke \& Behnke,
\href{https://doi.org/10.1007/978-3-531-90003-2_26}{2006}) and provides
significance level determinations for frequency parameters, thus serving
as an \emph{approximation} \emph{for exact} level determination through
probability distributions such as \emph{binomial}, \emph{multinomial} or
\emph{hypergeometric} (c.f. Camilli,
\href{https://doi.org/10.1007/BF02301418}{1995}), s. in this context
Wilson and Hilferty (\href{http://www.jstor.org/stable/86022}{1931}),
Patnaik (\href{http://www.jstor.org/stable/2332542}{1949}), Molinari
(\href{http://www.jstor.org/stable/2335780}{1977}), Hafner
(\href{https://doi.org/10.1007/978-3-7091-3420-7_14}{1992}), Jolicoeur
(\href{https://doi.org/10.1007/978-1-4615-4777-8_8}{1999}), Canal
(\href{https://doi.org/10.1016/j.csda.2004.04.001}{2005}), Brereton
(\href{https://doi.org/10.1002/cem.2734}{2015}) and Das
(\href{https://arxiv.org/abs/2404.05062}{2025}).

(1.14.-1.15.) \emph{Normal} distribution (de Moivre,
\href{https://books.google.com/books?id=PII_AAAAcAAJ}{1738}) or
\emph{Gaussian} distribution is fundamental to the \emph{central limit
theorem}, which states that the means of many independent, identically
distributed random variables tend to follow this \emph{normal}
distribution, the \emph{basis} of \emph{inferential statistics}. The
\(z\)-transformation with a mean of \(\overline{x\ } = \ 0\) and a
standard deviation of \(s = 1\) yields the \emph{standard} \emph{normal}
distribution, c.f. Behnke and Behnke
(\href{https://doi.org/10.1007/978-3-531-90003-2_26}{2006}). See
moreover David (\href{http://www.jstor.org/stable/2332676}{1949}),
Breitenberger (\href{http://www.jstor.org/stable/2333749}{1963}),
Mulholland and Jones
(\href{https://doi.org/10.1007/978-1-4899-6507-3_8}{1968b}), Chung
(\href{https://doi.org/10.1007/978-1-4757-3973-2_7}{1974}), Leaver and
Thomas (\href{https://doi.org/10.1007/978-1-349-01942-7_4}{1974}), Thome
(\href{http://www.jstor.org/stable/40986007}{1990}), Sachs
(\href{https://doi.org/10.1007/978-3-642-77717-2_4}{1993}), Massart et
al. (\href{https://doi.org/10.1016/S0922-3487(97)80033-0}{1998}), von
Storch and Zwiers
(\href{https://doi.org/10.1017/CBO9780511612336}{1999}), Stahl
(\href{http://www.jstor.org/stable/27642916}{2006}), Tóth
(\href{https://doi.org/10.1002/cem.1382}{2011}), Wang
(\href{https://doi.org/10.1016/j.spl.2011.03.037}{2011}), Brereton
(\href{https://doi.org/10.1002/cem.2655}{2014}), Lyon
(\href{https://doi.org/10.1093/bjps/axs046}{2014}), Bera et al.
(\href{http://www.jstor.org/stable/43948011}{2016}), Wesolowski and
Musselwhite Thompson
(\href{https://doi.org/10.4135/9781506326139.n476}{2018}), Benzon
(\href{https://doi.org/10.48188/so.2.6}{2021}) or Saneii and Doosti
(\href{https://doi.org/10.1007/978-981-97-3083-4_6}{2024}).

\emph{Theta} applications (1.16.-1.23.) generating distributions and
\emph{estimators} for several parameters \(\theta\) (\emph{theta})
within different designs via \emph{bootstrap} method (Efron,
\href{https://doi.org/10.1214/aos/1176344552}{1979},
\href{https://doi.org/10.1093/biomet/68.3.589}{1981},
\href{https://doi.org/10.1137/1.9781611970319}{1982}, res.), with given
number of resamples \(B\), where \emph{bootstrap} estimator

\[{\widehat{\theta}}_{B} = B^{- 1} \cdot \sum_{i = 1}^{B}\theta_{i}^{*}.\]

Implemented in \texttt{ThetaWin} (see Schrausser,
\href{https://www.academia.edu/81800920}{2009},
\href{https://doi.org/10.5281/zenodo.7659264}{2023g},
\href{https://doi.org/10.5281/zenodo.17880113}{2025a}); for a deeper and
current insight into the \emph{bootstrap} method see e.g.~Politis
(\href{https://doi.org/10.1109/79.647042}{1998}), Bickel and Ren
(\href{http://www.jstor.org/stable/4356107}{2001}), Horowitz
(\href{https://doi.org/10.1016/S1573-4412(01)05005-X}{2001}), Wilcox
(\href{https://doi.org/10.1007/978-1-4757-3522-2_6}{2001}), Machado and
Parente (\href{http://www.jstor.org/stable/23114968}{2005}), Kleiner et
al. (\href{http://www.jstor.org/stable/24774569}{2014}), Kenett et al.
(\href{https://doi.org/10.1007/978-3-031-07566-7_3}{2022}) or Zuev
(\href{https://doi.org/10.1007/978-3-032-03848-7_6}{2026}); furthermore
Hall (\href{http://www.jstor.org/stable/3182845}{2003}), Politis
(\href{https://doi.org/10.1214/ss/1063994977}{2003}), Pewsey
(\href{https://doi.org/10.2307/2348962}{2018}), Khosravi et al.
(\href{https://doi.org/10.1016/j.seps.2020.100781}{2021}), Baíllo and
Cárcamo (\href{https://doi.org/10.1007/s11222-025-10762-z}{2025}),
Eidous (\href{https://doi.org/10.13140/RG.2.2.12238.11846}{2025}), Liu
(\href{https://doi.org/10.1007/s41237-025-00283-4}{2025}) or Zheng and
Fan (\href{https://doi.org/10.3390/math13182913}{2025}).

\hypertarget{binomial}{%
\subsubsection{1.1. Binomial}\label{binomial}}

\begin{center}\rule{0.5\linewidth}{0.5pt}\end{center}

\[f\left( X \leq k|n \right) = \sum_{i = 0}^{k}\frac{n!}{i! \cdot (n - i)!} \cdot p^{i} \cdot q^{(n - i)}.\]

Usage:

\begin{verbatim}
Binomial [p][k][n] [[1]]
[p] ............. Probability of event A
[k] ............. n of events A
[n] ............. n of trials
[1] ............. (1):full output
\end{verbatim}

\hypertarget{binomial_t}{%
\subsubsection{1.2. Binomial\_T}\label{binomial_t}}

\begin{center}\rule{0.5\linewidth}{0.5pt}\end{center}

\[f\left( X \leq b|b,c \right) = p = \sum_{i = 0}^{b}\frac{(b + c)!}{i! \cdot (b + c - i)!} \cdot 2^{- i} \cdot 2^{- (b + c - i)}.\]

Usage:

\begin{verbatim}
Binomial_T [b][c] [[1]]
[b] ............. Cell count b
[c] ............. Cell count c
[1] ............. (1):full output
\end{verbatim}

\hypertarget{epsilon}{%
\subsubsection{1.3. Epsilon}\label{epsilon}}

\begin{center}\rule{0.5\linewidth}{0.5pt}\end{center}

\[\varepsilon = \frac{\mu_{1} - \mu_{0}}{\widehat{\sigma}}.\]

Usage:

\begin{verbatim}
epsilon [mode][Q0][s][n][e|Q1][p][df] [[x]]
[mode] .......... (1):Effect-size (2):Theta.1
[Q0] ............ Theta.0
[s] ............. Standard deviation
[n] ............. n of cases
[e|Q1] .......... Epsilon | Theta.1
[p] ............. Percent-level (0.00)
[df] ............ Degrees of freedom n - (.)
[x] ............. Test value
\end{verbatim}

\hypertarget{f_function}{%
\subsubsection{1.4. F\_Function}\label{f_function}}

\begin{center}\rule{0.5\linewidth}{0.5pt}\end{center}

\[F\left( F,df_{1},df_{2} \right) = 1 - p^{\alpha 2} = \int_{0}^{F}\frac{\Gamma_{\left( \frac{df_{1} + df_{2}}{2} \right)}}{\Gamma_{\left( \frac{df_{1}}{2} \right)} \cdot \Gamma_{\left( \frac{df_{2}}{2} \right)}} \cdot \left( \frac{df_{1}}{df_{2}} \right)^{\frac{df_{1}}{2}} \cdot F^{\frac{df_{1}}{2} - 1} \cdot \left( 1 + \frac{df_{1}}{df_{2}} \cdot F \right)^{- \frac{df_{1} + df_{2}}{2}}\ dF.\]

Usage:

\begin{verbatim}
F_Function [mode][x][n1][n2]
[mode] .......... (1):Fx=p->F (2):Fy=F->p
[x] ............. p-value/F-value
[n1] ............ n1
[n2] ............ n2
\end{verbatim}

\hypertarget{fisher_exact}{%
\subsubsection{1.5. Fisher\_Exact}\label{fisher_exact}}

\begin{center}\rule{0.5\linewidth}{0.5pt}\end{center}

\[f\left( X = a|a,b,c,d \right) = P0 = \frac{(a + b)! \cdot (c + d)! \cdot (a + c)! \cdot (b + d)!}{(a + b + c + d)! \cdot a! \cdot b! \cdot c! \cdot d!},\]

\[f\left( X \leq n|a,b,c,d \right) = p^{exact2} = \sum_{i = 0}^{n}Pi;\ Pi \leq P0.\]

Usage:

\begin{verbatim}
Fisher_Exact [a][b][c][d] [[1]]
[a][b][c][d] .... Cell counts a,b,c,d
[1] ............. (1):full output
\end{verbatim}

\hypertarget{fisher_z}{%
\subsubsection{1.6. Fisher\_Z}\label{fisher_z}}

\begin{center}\rule{0.5\linewidth}{0.5pt}\end{center}

\[Z = \frac{1}{2} \cdot log_{e}\left( \frac{1 + r}{1 - r} \right),\]
\[r = \frac{e^{2 \cdot Z} - 1}{e^{2 \cdot Z} + 1}.\]

Usage:

\begin{verbatim}
Fisher_Z [mode][x]
[mode] ......... (1):r->Z (2):Z->r
[x] ............ r-value/Z-value
\end{verbatim}

\hypertarget{gamma_function}{%
\subsubsection{1.7. GAMMA\_Function}\label{gamma_function}}

\begin{center}\rule{0.5\linewidth}{0.5pt}\end{center}

\[f(x,t) = \Gamma = \int_{0}^{\infty}t^{x - 1} \cdot e^{- t}\ dt + c.\]

Usage:

\begin{verbatim}
GAMMA_Function [mode][value]
[mode] ......... (1):F(x)->⌈ (2):F'(⌈)->x
[value] ........ x / ⌈
\end{verbatim}

\hypertarget{gamma}{%
\subsubsection{1.8. GAMMA}\label{gamma}}

\begin{center}\rule{0.5\linewidth}{0.5pt}\end{center}

Usage:

\begin{verbatim}
GAMMA [n][input][output]
[n] ............ n of cases
[input] ........ Input file
[output] ....... Output file
\end{verbatim}

\hypertarget{geometric}{%
\subsubsection{1.9. Geometric}\label{geometric}}

\begin{center}\rule{0.5\linewidth}{0.5pt}\end{center}

\[f\left( X \leq r|p \right) = \sum_{i = 0}^{r}{p \cdot q^{i}}.\]

Usage:

\begin{verbatim}
Geometric [p][r+1] [[1]]
[p] ............ Probability of event A
[r+1] .......... n of trials
[1] ............ (1):full output
\end{verbatim}

\hypertarget{hypergeometric}{%
\subsubsection{1.10. Hypergeometric}\label{hypergeometric}}

\begin{center}\rule{0.5\linewidth}{0.5pt}\end{center}

\[f\left( X \leq k|n,K,N \right) = \sum_{i = 0}^{k}\frac{\binom{K}{i} \cdot \binom{N - K}{n - i}}{\binom{N}{n}}.\]

Usage:

\begin{verbatim}
Hypergeometric [k][n][N][K] [[1]]
[k] ............ n of events A in Sub-Population
[n] ............ Size of Sub-Population
[N] ............ Size of Population
[K] ............ n of events A in Population
[1] ............ (1):full output
\end{verbatim}

\hypertarget{poisson}{%
\subsubsection{1.11. Poisson}\label{poisson}}

\begin{center}\rule{0.5\linewidth}{0.5pt}\end{center}

\[f\left( X \leq k|n,p \right) = \sum_{i = 0}^{k}\frac{(n \cdot p)^{i}}{e^{n \cdot p} \cdot i!}.\]

Usage:

\begin{verbatim}
Poisson [p][k][n] [[1]]
[p] ............ Probability of event A
[k] ............ n of events A
[n] ............ n of trials
[1] ............ (1):full output
\end{verbatim}

\hypertarget{t_function}{%
\subsubsection{1.12. t\_Function}\label{t_function}}

\begin{center}\rule{0.5\linewidth}{0.5pt}\end{center}

\[F(t,df) = p = \int_{- \infty}^{t}\frac{\Gamma_{\left( \frac{df - 1}{2} \right)}}{\Gamma_{\left( \frac{df}{2} \right)}} \cdot (df \cdot \pi)^{- \frac{1}{2}} \cdot \left( 1 + \frac{t^{2}}{df} \right)^{- \frac{df - 1}{2}}\ dt.\]

Usage:

\begin{verbatim}
t_Function [mode][x][n]
[mode] ......... (1):Fx=p->t (2):Fy=t->p
[x] ............ p-value/t-value
[n] ............ n of cases
\end{verbatim}

\hypertarget{x2_function}{%
\subsubsection{1.13. x2\_Function}\label{x2_function}}

\begin{center}\rule{0.5\linewidth}{0.5pt}\end{center}

\[F\left( \chi^{2},df \right) = 1 - p^{\alpha 2} = \int_{0}^{\chi^{2}}\frac{1}{2^{\frac{df}{2}} \cdot \Gamma_{\left( \frac{df}{2} \right)}} \cdot \left( \chi^{2} \right)^{\frac{df}{2} - 1} \cdot e^{- \frac{\chi^{2}}{2}}\ d\chi^{2}.\]

Usage:

\begin{verbatim}
x2_Function [mode][x][n]
[mode] ......... (1):Fx=p->x² (2):Fy=x²->p
[x] ............ p-value/x²-value
[n] ............ n of cases
\end{verbatim}

\hypertarget{z_dichte}{%
\subsubsection{1.14. z\_Dichte}\label{z_dichte}}

\begin{center}\rule{0.5\linewidth}{0.5pt}\end{center}

\[f(z) = \vartheta = \frac{1}{\sqrt{2 \cdot \pi}} \cdot e^{- \frac{z^{2}}{2}},\]
\[f^{- 1}(z) = f(\vartheta) = z = \sqrt{ln(\frac{\vartheta}{\sqrt{(2 \cdot \pi)^{- 1}}})^{- 2}},\]
\[F(z) = p = \int_{- \infty}^{z}\vartheta\ dz.\]

Usage:

\begin{verbatim}
z_Dichte [mode][value] [[f]]
[mode] ......... (1):fx=z->d (2):fy=d->z (3):∫x=z->p
                 (4):∫'p->z (5):∫'p->d (6):∫y=d->p
[value] ........ z-value/z-density/percent-rank p
[f] ............ (1):z-density function graph
\end{verbatim}

\hypertarget{z_function}{%
\subsubsection{1.15. z\_Function}\label{z_function}}

\begin{center}\rule{0.5\linewidth}{0.5pt}\end{center}

Usage:

\begin{verbatim}
z_Function [mode][x]
[mode] ......... (1):Fx=p->z (2):Fy=z->p
[x] ............ p-value/z-value
\end{verbatim}

\hypertarget{theta}{%
\subsubsection{1.16. Theta}\label{theta}}

\begin{center}\rule{0.5\linewidth}{0.5pt}\end{center}

Usage:

\begin{verbatim}
Theta [sd][min][max][qq][q][v][s] [[x]] [[g]]
[sd]  .......... Seed: |0| Timevalue
[min] .......... R Minimum value
[max] .......... R Maximum value
[qq]  .......... Theta-Theta/
[q]   .......... Theta:
                 |0| Harmonic mean (HM)
                 |1| Arithmetic mean (AM)
                 |2| Sum (SUM)
                 |3| Standard deviation (SD)
                 |4| Population variance estimation (VAR)
                 |5| Product sum (PSM)
                 |6| Geometric mean (GM)
                 |7| Schrausser's d (D)
                 |8| DvarO (DV)
[v]  ..........  n of Theta (v)
[s]  ..........  n Subpopulations (s)
[x]  ..........  Test value x
[g]  ..........  |1| Value range integer
\end{verbatim}

\hypertarget{theta_q}{%
\subsubsection{1.17. Theta\_Q}\label{theta_q}}

\begin{center}\rule{0.5\linewidth}{0.5pt}\end{center}

Usage:

\begin{verbatim}
Theta_Q [sd][min][max][qq][qp][qs1][qs2][qQ][v][m][n][s] [[x]] [[g]]
[sd]  ......................... Seed: |0| Timevalue
[min] ......................... R Minimum value
[max] ......................... R Maximum value
[qq]  ......................... Theta-Theta/
[qp]  ......................... Theta P/
[qs1] [qs2] ................... Theta S1, S2:
                                |0| Harmonic mean (HM)
                                |1| Arithmetic mean (AM)
                                |2| Sum (SUM)
                                |3| Standard deviation (SD)
                                |4| Population variance estm. (VAR)
                                |5| Product sum (PSM)
                                |6| Geometric mean (GM)
                                |7| Schrausser's d (D)
                                |8| DvarO (DV)
[qQ]  ......................... Theta Q:
                                |1| Difference
                                |2| Quotient
                                |3| Sum
                                |4| Product
[v]  .........................  n of Theta P (v)
[m]  .........................  n of Theta S1 (m)
[n]  .........................  n of Theta S2 (n)
[s]  .........................  n Subpopulations (s)
[x]  .........................  Test value x
[g]  .........................  |1| Value range integer
\end{verbatim}

\hypertarget{theta_qv}{%
\subsubsection{1.18. Theta\_Qv}\label{theta_qv}}

\begin{center}\rule{0.5\linewidth}{0.5pt}\end{center}

Usage:

\begin{verbatim}
Theta_Qv [sd][min][max][qq][qp][qs1][qs2][qQ][QQ][v][n][s] [[x]] [[g]]
[sd]  ......................... Seed: |0| Timevalue
[min] ......................... R Minimum value
[max] ......................... R Maximum value
[qq]  ......................... Theta-Theta/
[qp]  ......................... Theta P/
[qs1][qs2]..................... Theta S1, S2/
[qQ]  ......................... Theta Q:
                                |0| Harmonic mean (HM)
                                |1| Arithmetic mean (AM)
                                |2| Sum (SUM)
                                |3| Standard deviation (SD)
                                |4| Population variance estm. (VAR)
                                |5| Product sum (PSM)
                                |6| Geometric mean (GM)
                                |7| Schrausser's d (D)
                                |8| DvarO (DV)
[QQ]  ......................... Theta Theta Q:
                                |1| Difference
                                |2| Quotient
                                |3| Sum
                                |4| Product
                                |5| Correlation
                                |6| Covariance
                                |7| Coefficient of determination
                                |8| Redundancy
[v]  .......................... n of Theta P (v)
[n]  .......................... n of Theta S1,S2 (n)
[s]  .......................... n Subpopulations (s)
[x]  .......................... Test value x
[g]  .......................... |1| Value range integer
\end{verbatim}

\hypertarget{theta_rq}{%
\subsubsection{1.19. Theta\_rQ}\label{theta_rq}}

\begin{center}\rule{0.5\linewidth}{0.5pt}\end{center}

Usage:

\begin{verbatim}
Theta_rQ [sd][min][max][qq][qp][q11][q12][q21][q22][qr1][qr2][qQ][v]
[m][n][s] [[x]] [[g]]
[sd]  ......................... Seed: |0| Timevalue
[min] ......................... R Minimum value
[max] ......................... R Maximum value
[qq]  ......................... Theta-Theta/
[qp]  ......................... Theta P/
[q11][q12] .................... Theta S11, S12/
[q21][q22] .................... Theta S21, S22:
                                |0| Harmonic mean (HM)
                                |1| Arithmetic mean (AM)
                                |2| Sum (SUM)
                                |3| Standard deviation (SD)
                                |4| Population variance estm. (VAR)
                                |5| Product sum (PSM)
                                |6| Geometric mean (GM)
                                |7| Schrausser's d (D)
                                |8| DvarO (DV)
[qr1][qr2] .................... Theta Regressions 1,2/
                                |1| Correlation (kor)
                                |2| Covariance (cov)
                                |3| Coeff. of determination (det)
                                |4| Redundancy (red)
[qQ]  ......................... Theta Q:
                                |1| Difference (Diff)
                                |2| Quotient (Quot)
                                |3| Sum (Summ)
                                |4| Product (Prod)
[v]  .......................... n of Theta P (v)
[m]  .......................... n of Theta S11,S12 (m)
[n]  .......................... n of Theta S21,S22 (n)
[s]  .......................... n Subpopulations (s)
[x]  .......................... Test value x
[g]  .......................... |1| Value range integer
\end{verbatim}

\hypertarget{theta_s}{%
\subsubsection{1.20. Theta\_S}\label{theta_s}}

\begin{center}\rule{0.5\linewidth}{0.5pt}\end{center}

Usage:

\begin{verbatim}
Theta_S [sd][min][max][qq][qp][qs][v][m][s] [[x]] [[g]]
[sd]  ......................... Seed: |0| Timevalue
[min] ......................... R Minimum value
[max] ......................... R Maximum value
[qq]  ......................... Theta-Theta:
[qp]  ......................... Theta P/
[qs]  ......................... Theta S/
                                |0| Harmonic mean (HM)
                                |1| Arithmetic mean (AM)
                                |2| Sum (SUM)
                                |3| Standard deviation (SD)
                                |4| Population variance estm. (VAR)
                                |5| Product sum (PSM)
                                |6| Geometric mean (GM)
                                |7| Schrausser's d (D)
                                |8| DvarO (DV)
 [v]  ......................... n of Theta P (v)
 [m]  ......................... n of Theta S (m)
 [s]  ......................... n Subpopulations (s)
 [x]  ......................... Test value x
 [g]  ......................... |1| Value range integer
\end{verbatim}

\hypertarget{verteilungsform}{%
\subsubsection{1.21. Verteilungsform}\label{verteilungsform}}

\begin{center}\rule{0.5\linewidth}{0.5pt}\end{center}

Usage:

\begin{verbatim}
Verteilungsform [min][max][n][s]
[min] ......................... Minimum value
[max] ......................... Maximum value
[n] ............. n of parameter Theta=sum(x)
[s] ........................ n Subpopulations
\end{verbatim}

\hypertarget{verteilungsform_2u}{%
\subsubsection{1.22. Verteilungsform\_2u}\label{verteilungsform_2u}}

\begin{center}\rule{0.5\linewidth}{0.5pt}\end{center}

Usage:

\begin{verbatim}
Verteilungsform_2u [min][max][q][n1][n2][s] [[xd]] [[g]]
[min] ......................... Minimum value
[max] ......................... Maximum value
[q] ........................... Theta:
                                |0| Harmonic mean
                                |1| Arithmetic mean
                                |2| Sum
                                |3| Standard deviation 
                                |4| Population variance estm.
                                |5| Product sum 
                                |6| Geometric mean
                                |7| Schrausser's d
                                |8| DvarO
[n1] .......................... n1 of Theta
[n2] .......................... n2 of Theta
[s] ........................... n Subpopulations
[xd] .......................... Difference test value
[g] ........................... |1| Value range integer
\end{verbatim}

\hypertarget{verteilungsform_kor}{%
\subsubsection{1.23. Verteilungsform\_kor}\label{verteilungsform_kor}}

\begin{center}\rule{0.5\linewidth}{0.5pt}\end{center}

Usage:

\begin{verbatim}
Verteilungsform_kor [min][max][q][n][s] [[x]] [[g]]
[min] ......................... Minimum value
[max] ......................... Maximum value
[q] ........................... Theta:
                                |1| Pearson Correlation
                                |2| Covariance
                                |3| Coeff. of determination
                                |4| Redundancy
                                |5| Regression coefficient ayx
                                |6| Regression coefficient byx
                                |7| Regression coefficient axy
                                |8| Regression coefficient bxy
[n] ........................... n of Theta
[s] ........................... n Subpopulations
[x] ........................... Test value x
[g] ........................... |1| Value range integer
\end{verbatim}

\textbf{Table} 1. \emph{ConsoleApp} \emph{distribution} functions and
\emph{bootstrap} overview with corresponding functions in
\texttt{SCHRAUSSER-MAT} and \texttt{HP\_Prime\_MATH} (Schrausser,
\href{https://doi.org/10.17605/OSF.IO/8XE42}{2022},
\href{https://doi.org/10.5281/zenodo.15713317}{2025a}, res.)\emph{.}

\begin{verbatim}
nr.     Function        
        ConsoleApp              SCHRAUSSER-MAT                  HP_Prime_MATH

1       Binomial¹               BNP, BNW    
2       Binomial_T¹             BN1, BN2                        ABT1, BINOM, pzBN, zBN, E01
3       Epsilon¹                EFG, ANI, BNI, IMB, EFS, OPP    EPSILON, EFG, EFR, E01
4       F_Function¹             FPW, PFW, PFD                   FVTLG, F02, F03Z
5       Fisher_Exact¹           FX0, FX1, FX2                   FX_, z4F, pz4F
6       Fisher_Z¹               FZR, RFZ                        Zcor, rZ, Zr
7       GAMMA_Function¹         GAMMA, AGAM 
8       GAMMA¹                  IGM F01z, F04
9       Geometric¹              GMP, GMW                        GMVTLG
10      Hypergeometric¹         HGP, HGW    
11      Poisson¹                PNP, PNW    
12      t_Function¹             TPW, PTW, PTD                   tVTLG, F06_, F02, F06, F03Z
13      x2_Function¹            XPW, PXW, PXD                   ch2VTLG, F07_, F02, F07, F03Z
14      z_Dichte¹               DZW, ZDW    
15      z_Function¹             ZPW, PZW, PZD                   NVTLG, E01, F02, F03
16      Theta²      
17      Theta_Q²        
18      Theta_Qv²       
19      Theta_rQ²       
20      Theta_S²        
21      Verteilungsform     
22      Verteilungsform_2u      
23      Verteilungsform_kor
\end{verbatim}

¹) implemented in \texttt{FunktionWin} (Schrausser,
\href{https://doi.org/10.5281/zenodo.7651661}{2023f},
\href{https://doi.org/10.5281/zenodo.17880113}{2025c}).\\
²) implemented in \texttt{ThetaWin} (Schrausser,
\href{https://doi.org/10.5281/zenodo.7659264}{2023g},
\href{https://doi.org/10.5281/zenodo.17880113}{2025a}).

\hypertarget{consoleapp_eeg}{%
\subsection{2. ConsoleApp\_EEG}\label{consoleapp_eeg}}

Console applications (c.f. Schrausser,
\href{https://doi.org/10.5281/zenodo.10701349}{2024b}) for \emph{EEG}
(s. Galvani,
\href{https://doi.org/10.5479/sil.324681.39088000932442}{1791}; Caton,
\href{https://journals.lww.com/jonmd/fulltext/1875/10000/electrical_currents_of_the_brain.13.aspx}{1875};
Beck,
\href{https://scholar.google.com/scholar_lookup?\&title=Die\%20Str\%C3\%B6me\%20der\%20Nervenzentren\&journal=Zbl\%20Physiol\&volume=4\&pages=572-573\&publication_year=1890\&author=Beck\%2CA}{1890}
and Berger, \href{https://doi.org/10.1007/BF01797193}{1929}) parameters
implemented in \texttt{ATINA} (Schrausser,
\href{https://doi.org/10.5281/zenodo.19422127}{2006}), calculating
\emph{coherence} (2.1.; c.f. French \& Graham,
\href{https://doi.org/10.1016/0167-8760(84)90044-8}{1984}; Schrausser et
al., \href{https://doi.org/10.5281/zenodo.13738772}{2001}; Guevara et
al., \href{https://doi.org/10.4236/jbise.2011.412096}{2011}; Puthanmadam
Subramaniyam \& Thiagarajan,
\href{https://doi.org/10.1038/s41598-025-94076-0}{2025}), \emph{cross}
\emph{correlation} (2.2.) and \emph{focus} parameters (2.3.-2.6.; s.
Schrausser, \href{https://doi.org/10.13140/RG.2.2.32114.17601}{2000a},
\href{http://doi.org/10.13140/RG.2.2.28637.90083}{b}) from
\emph{Fourier-analyzed} data (c.f. Fourier,
\href{https://archive.org/details/bub_gb_TDQJAAAAIAAJ/mode/1up}{1822}).
More profound insight into the methods is provided by e.g.~Gerthsen
(\href{https://doi.org/10.1007/978-3-662-30201-9}{1966}), Kumar and
Bhuvaneswari
(\href{https://doi.org/10.1016/j.proeng.2012.06.298}{2012}), Siuly et
al.~(\href{https://doi.org/10.1007/978-3-319-47653-7}{2017}), Panov
(\href{https://doi.org/10.1007/978-1-0716-3230-7_5}{2024}), Grafakos
(\href{https://doi.org/10.1007/978-3-031-56500-7}{2024}), Panitz et al.
(\href{https://doi.org/10.1017/9781009000796.024}{2024}), Brigola
(\href{https://doi.org/10.1007/978-3-031-81311-5}{2025}) or Meschede et
al. (\href{https://doi.org/10.1007/978-3-662-30201-9}{2026}).

For the implementation of the methods in the fields of neuropsychology
or psychobiology see e.g.~Neubauer et al.
(\href{https://doi.org/10.1016/S0160-2896(02)00091-0}{2002}), Neubauer
et al. (\href{https://doi.org/10.1016/j.cogbrainres.2005.05.011}{2005}),
Micheloyannis et al.
(\href{https://doi.org/10.1016/j.neulet.2006.04.006}{2006}), Neubauer
and Fink (\href{https://doi.org/10.1016/j.neubiorev.2009.04.001}{2009}),
Chiarionet
al.~(\href{https://doi.org/10.3390/bioengineering10030372}{2023}), Zhang
et al.~(\href{https://doi.org/10.1186/s40779-023-00502-7}{2023}) or
Nguyen et al.
(\href{https://doi.org/10.1161/circ.152.suppl/_3.4373687}{2025}), c.f.
further Adrian (\href{https://doi.org/10.9783/9781512809794}{2016}),
Tassinari (\href{https://doi.org/10.1016/j.clinph.2019.10.010}{2019})
and Rossini et
al.~(\href{https://doi.org/10.1016/j.clinph.2024.11.021}{2025}).

\hypertarget{erc-ercx}{%
\subsubsection{2.1. ERC, ERCX}\label{erc-ercx}}

\begin{center}\rule{0.5\linewidth}{0.5pt}\end{center}

Calculates \emph{event related} \emph{coherence} \(ERC\), where
\emph{coherence}

\[Coh_{xy}^{2} = \frac{P_{xy}^{2}(f)}{P_{xx}(f) \cdot P_{yy}(f)},\]

with \emph{cross power} \(P_{xy}\) within given frequency \(f\) in
\(\mathbb{C}\) defined by

\[P_{xy}(f)\mathfrak{= R}\left( a_{xy} \right)^{2}\mathfrak{+ I}\left( a_{xy} \right)^{2},\]

c.f. Schrausser
(\href{http://doi.org/10.13140/RG.2.2.28637.90083}{2000b}).

Usage:

\begin{verbatim}
erc [input][output][nCoh][nX][refs][refe][acts][acte][typ]
[input] ................. Input File, Format ASCII tab. (e.g. coh.asc)
[output] ................ Output File
[nCOH] .................. Number of coherence values 
[nX] .................... Number of channel combinations
[refs] .................. Number of reference start value 
[refe] .................. Number of reference end value
[acts] .................. Number of active start value
[acte] .................. Number of active end value
[typ] ................... Type of output (0):individual (1):append
\end{verbatim}

\hypertarget{xcor}{%
\subsubsection{2.2. XCOR}\label{xcor}}

\begin{center}\rule{0.5\linewidth}{0.5pt}\end{center}

Calculates \emph{cross correlation} \(xCOR\).

Usage:

\begin{verbatim}
xcor [input][output][n][k]
[input] ..... Input File, Format ASCII tab. (e.g. data.dat)
[output] .... Output File
[n] ......... Number of cases
[k] ......... Number of variables
\end{verbatim}

\hypertarget{foc}{%
\subsubsection{2.3. FOC}\label{foc}}

\begin{center}\rule{0.5\linewidth}{0.5pt}\end{center}

Calculates \emph{focus} parameter \(yf\), where

\[yf = \sum_{i = 1}^{k}{1 - \frac{\frac{x_{i} - x_{\min}}{d}}{k - 1}},\]

with

\[d = x_{\max} - x_{\min}.\]

Usage:

\begin{verbatim}
foc [input][output][n][k]
[input] ..... Input File, Format ASCII tab. (e.g. data.dat)
[output] .... Output File
[n] ......... Number of cases
[k] ......... Number of leads
\end{verbatim}

\hypertarget{dis}{%
\subsubsection{2.4. DIS}\label{dis}}

\begin{center}\rule{0.5\linewidth}{0.5pt}\end{center}

Reading in \emph{distance} values \(d_{ij}\) for \(x_{i}\), \(x_{j}\)
and calculates \emph{wights} \(g_{ij}\) defined by

\[g_{ij} = \frac{\frac{1}{d_{ij}}}{\sum_{j = 1}^{n}\frac{1}{d_{ij}}}.\]

Usage:

\begin{verbatim}
dis [input][output][n][k]
[input] ..... Input File, Format ASCII tab. 
[output] .... Output File
[n] ......... Number of cases
[k] ......... Number of leads
\end{verbatim}

\hypertarget{floc}{%
\subsubsection{2.5. FLOC}\label{floc}}

\begin{center}\rule{0.5\linewidth}{0.5pt}\end{center}

Calculates \emph{spatial} \emph{focus} parameter \(yloc\), where

\[yloc = \frac{\sum_{i = 1}^{k}{\sum_{j = 1}^{k - 1}{\frac{1}{G_{i}} \cdot \left\lbrack \left( 1 - \frac{x_{i} - x_{\min}}{d} \right) \cdot w_{G_{ij}} \right\rbrack}}}{k - 2},\]

with

\[G_{i} = \sum_{j = 1}^{k - 1}g_{ij},\]
\[w_{G_{ij}} = g_{ij} \cdot \left( 1 - \frac{x_{j} - x_{\min}}{d} \right).\]

Usage:

\begin{verbatim}
floc [input][output][n][k][distance]
[input] ..... Input File, Format ASCII tab. (e.g. data.dat)
[output] .... Output File
[n] ......... Number of cases
[k] ......... Number of leads
[distance] .. Distance Matrix File
\end{verbatim}

\hypertarget{out}{%
\subsubsection{2.6. OUT}\label{out}}

\begin{center}\rule{0.5\linewidth}{0.5pt}\end{center}

Generates output file.

Usage:

\begin{verbatim}
out [input][n][k][start][end] 
[input] ..... Input File, Format ASCII tab. (e.g. coh.asc)
[n] ......... Number of values / cases / rows
[k] ......... Number of variables / columns 
[start] ..... Number of block start value 
[end] ....... Number of block end value
\end{verbatim}

\hypertarget{consoleapp_integral}{%
\subsection{3. ConsoleApp\_Integral}\label{consoleapp_integral}}

Console applications for \emph{integral} and \emph{interpolation} (see
Schrausser, \href{https://doi.org/10.5281/zenodo.7655056}{2023b}). For
the basic concept of \emph{Riemann sums} (Riemann,
\href{http://eudml.org/doc/135759}{1868}), see Lewis et al.
(\href{https://doi.org/10.1007/978-3-0348-5563-1_25.}{1978}), Hassler
(\href{https://doi.org/10.1007/978-3-540-73568-7_7}{2007}), Truc
(\href{https://www.jstor.org/stable/48662056}{2019}), Torchinsky
(\href{https://doi.org/10.1007/978-3-031-11799-2}{2022}), Forster and
Lindemann (\href{https://doi.org/10.1007/978-3-658-40130-6_18}{2023}),
further Goel and Rodriguez
(\href{http://www.jstor.org/stable/2689344}{1987}), Guillemin and
Stroock (\href{https://doi.org/10.1007/978-0-8176-4683-7_1}{2008}),
Büchter and Henn
(\href{https://doi.org/10.1007/978-3-8274-2680-2_6}{2010}), for
\emph{integral calculus} in general see Schwarz
(\href{https://doi.org/10.1007/978-3-663-01227-6_8}{1997}), Meyberg and
Vachenauer (\href{https://doi.org/10.1007/978-3-642-56654-7_4}{2001a}),
Allen and Isaacson
(\href{https://doi.org/10.1002/9781119245476.ch7}{2019}) or Neher
(\href{https://doi.org/10.1007/978-3-662-68815-1_7}{2024}).

The included methods are as follows:

(3.1.-3.2.) \emph{Romberg's method} (Romberg,
\href{https://www.google.com/url?sa=t\&source=web\&rct=j\&opi=89978449\&url=https://web.eng.fiu.edu/LEVY/images/EGM5346/romberg\%2520textbook\%2520example.pdf\&ved=2ahUKEwiSxr-u846SAxVFQvEDHd83Bj8QFnoECHMQAQ\&usg=AOvVaw3vtOkKCF8WeDkOB49i_03i}{1955})
as an analytical method for achieving higher approximation accuracy for
integrals, based on the \emph{Richardson extrapolation} (Richardson,
\href{https://doi.org/10.1098/rsta.1911.0009}{1911}; s. Rannacher,
\href{https://doi.org/10.1007/978-3-322-85997-6_9}{1987}; Bickel \&
Yahav, \href{http://www.jstor.org/stable/2288854}{1988}; Sidi,
\href{https://doi.org/10.1007/978-3-0348-6398-8_22}{1988},
\href{https://doi.org/10.1017/CBO9780511546815}{2003}; Dubeau,
\href{https://doi.org/10.1016/j.jcpx.2019.100017}{2019}; Bach,
\href{https://doi.org/10.1137/21M1397349}{2021}) or the
\emph{trapezoidal rule} (c.f. Luke et al.,
\href{https://doi.org/10.1016/S0076-5392(09)60075-8}{1969}; Abdulhameed
\& Memon, \href{https://doi.org/10.1088/1742-6596/2090/1/012104}{2021};
Izzo et al., \href{https://doi.org/10.1016/j.jcp.2022.111193}{2022};
Bujang, \href{https://doi.org/10.1063/5.0294833}{2025}), respectively.
The \emph{Richardson extrapolation} is further used, for example, in
machine learning and other fields to \emph{accelerate} convergences. In
this context c.f. Herrmann
(\href{https://doi.org/10.1007/978-3-322-96321-5_27}{1983}), Jetter
(\href{https://doi.org/10.1007/BF01389471}{1984}), Talay and Tubaro
(\href{https://doi.org/10.1016/0167-4730(90)90036-O}{1990}), von
Petersdorff (\href{http://www.jstor.org/stable/2324787}{1993}) or
Meyberg and Vachenauer
(\href{https://doi.org/10.1007/978-3-642-56654-7_4}{2001b}, p.~209).

(3.3.-3.4.) the \emph{Cubic spline} interpolation (Schoenberg,
\href{https://doi.org/10.1090/qam/15914}{1946}; c.f. Rabut,
\href{https://doi.org/10.1016/0898-1221(92)90177-J}{1992}; Agarwal \&
Wong, \href{https://doi.org/10.1007/978-94-011-2026-5_6}{1993}; Farin,
\href{https://doi.org/10.1007/978-3-663-10602-9_9}{1994}) as a method
for \emph{polynomial} \emph{curve fitting} and precise interpolation of
\emph{empirical} data points, where the rate of change corresponds to
the slope, resulting in smooth, continuous function curves. See in this
context Cheng and Barsky
(\href{https://doi.org/10.1016/0010-4485(91)90023-P}{1991}), Fredenhagen
et al.~(\href{https://doi.org/10.1006/jath.1998.3247}{1999}), Jiang and
Zhao (\href{https://doi.org/10.1109/FSKD.2009.591}{2009}), Moon
(\href{https://doi.org/10.1016/S0096-3003(99)00178-2}{2001}), Sarfraz
(\href{https://doi.org/10.1007/978-1-84628-871-5_8}{2008}), Sun et al.
(\href{https://doi.org/10.1016/j.cam.2022.115039}{2023}), Török et al.
(\href{https://doi.org/10.1016/j.amc.2025.129411}{2025}), Jarre
(\href{https://doi.org/10.48550/arXiv.2507.05083}{2025},
\href{https://doi.org/10.1016/j.cam.2025.117240}{2026}). For application
in \emph{EEG} examination and psychophysiology, see e.g.~Perrin et al.
(\href{https://doi.org/10.1016/0013-4694(87)90141-6}{1987}), Soufflet et
al. (\href{https://doi.org/10.1016/0013-4694(91)90204-H}{1991}), Congedo
et al. (\href{https://doi.org/10.1300/J184v06n04_08}{2002}), Roy and
Shukla (\href{https://doi.org/10.1007/s11277-017-4846-3}{2017}),
Abduganiev et al.
(\href{https://doi.org/10.1007/978-3-031-27199-1_3}{2023}) or Schulter
et al. (\href{https://doi.org/10.1038/s41598-023-49307-7}{2023}, p.~5).

(3.5.-3.6.) the \emph{Newton} interpolation (Cardano,
\href{https://web.archive.org/web/20220201093634/http:/www.filosofia.unimi.it/cardano/testi/operaomnia/vol_4_s_4.pdf}{1545};
Newton,
\href{https://books.google.at/books?id=XJwx0lnKvOgC\&pg=PP2\&redir_esc=y\#v=onepage\&q\&f=false}{1687})
as a further method for \emph{curve fitting} which is easier to
calculate, however, has lower accuracy and stability compared to
interpolation via \emph{cubic splines}, calculates a significantly less
smooth curve and is therefore better suited for \emph{smaller} datasets
and lower accuracy requirements. For an overview c.f. Rutishauser
(\href{https://doi.org/10.1007/978-1-4612-3468-5_6}{1990}), Scherer
(\href{https://doi.org/10.1007/978-3-642-13990-1_2}{2010}) or Epperson
(\href{https://doi.org/10.1002/9781119604570.ch4}{2021}), further
insight give e.g.~Tsao (\href{https://doi.org/10.1007/BF01389317}{1977})
and Zou et al.~(\href{https://doi.org/10.1155/2020/9020541}{2020}).

\hypertarget{romi}{%
\subsubsection{3.1. ROMI}\label{romi}}

\begin{center}\rule{0.5\linewidth}{0.5pt}\end{center}

Approximates (i) \emph{surface} or (ii) \emph{curve integrals
{[}Approximiert (i) Flächen- oder (ii) Kurven-Integrale{]}}

\[\int_{a}^{b}f(x)dx\]

by means of the \emph{Romberg method}, if applicable including file
output to \textbf{\emph{romi.txt}} from the function \emph{matrix
{[}mittels Romberg Methode, dabei ggf. Dateiausgabe nach romi.txt von
Funktionsmatrix{]}}

\[\mathbf{F =}\left( \mathbf{x|y} \right),\]

where \emph{{[}bei{]}}

\[x = \frac{x + d}{2},y = \int_{a}^{b}{f(x)dx}.\]

\begin{itemize}
\item
  Execution of \textbf{\emph{ROMI.bat}}:
\item
  Definition of \(f(x)\) in \textbf{\emph{ROMI.h}};
\item
  Compiling \textbf{\emph{ROMI.c}};
\item
  Execution of \textbf{\emph{ROMI.exe}}.
\item
  \emph{{[}Ausführung von ROMI.bat:{]}}
\item
  \emph{{[}Definition von} \(f(x)\) \emph{in ROMI.h;{]}}
\item
  \emph{{[}Compilieren von ROMI.c;{]}}
\item
  \emph{{[}Ausführung von ROMI.exe.{]}}
\end{itemize}

Usage \emph{{[}Handhabung{]}}:

\begin{verbatim}
ROMI [a][b][d][m][F]
[a] ......... Integration Minimum a
[b] ......... Integration Maximum b
[d] ......... Delta d
[m] ......... Mode (0):Area.I (1):Curve.I
[F] ......... Function matrix (0):none (1):romi.txt
\end{verbatim}

Example \emph{{[}Bsp.{]}}:

\begin{verbatim}
ROMI.bat
ROMI 1 10 0.0001 0 0
ROMI -9 0 0.5 0 1
\end{verbatim}

\hypertarget{rome}{%
\subsubsection{3.2. ROME}\label{rome}}

\begin{center}\rule{0.5\linewidth}{0.5pt}\end{center}

Approximates \(\int_{a}^{b}f(x)dx\) using \emph{Romberg extrapolation}
\emph{{[}Approximiert} \(\int_{a}^{b}f(x)dx\) \emph{mittels
Romberg-Extrapolation{]}} (see. Meyberg \& Vachenauer,
\href{https://doi.org/10.1007/978-3-642-56654-7_4}{2001b}, p.~209).

\begin{itemize}
\item
  Execution of \textbf{\emph{ROME.bat}}:
\item
  Definition of \(f(x)\) in \textbf{\emph{ROME.h}}
\item
  Compiling \textbf{\emph{ROME.c}};
\item
  Execution of \textbf{\emph{ROME.exe}}.
\item
  \emph{{[}Ausführung von ROME.bat:{]}}
\item
  \emph{{[}Definition von} \(f(x)\) \emph{in ROME.h;{]}}
\item
  \emph{{[}Compilieren von ROME.c;{]}}
\item
  \emph{{[}Ausführung von ROME.exe.{]}}
\end{itemize}

Usage \emph{{[}Handhabung{]}}:

\begin{verbatim}
ROME [a][b]
[a] ......... Integration Minimum a
[b] ......... Integration Maximum b
\end{verbatim}

Example \emph{{[}Bsp.{]}}:

\begin{verbatim}
ROME.bat
ROME 1 10
ROME -14 1
\end{verbatim}

\hypertarget{kusi}{%
\subsubsection{3.3. KUSI}\label{kusi}}

\begin{center}\rule{0.5\linewidth}{0.5pt}\end{center}

\emph{Cubic spline} interpolation: Calculation of the \emph{coefficient}
\emph{matrix}
\(\mathbf{A =}\left( \mathbf{b}\left| \mathbf{c} \right|\mathbf{d} \right)\)
and \(\mathbf{s}\left( \mathbf{x} \right)\) for an (empirical) function
\emph{matrix} \(\mathbf{F =}\left( \mathbf{x|y} \right)\), where
\emph{{[}Kubische Spline Interpolation: Berechnung der
Koeffizientenmatrix}
\(\mathbf{A =}\left( \mathbf{b}\left| \mathbf{c} \right|\mathbf{d} \right)\)
\emph{sowie} \(\mathbf{s}\left( \mathbf{x} \right)\) \emph{zu einer
(empirischen) Funktionsmatrix}
\(\mathbf{F =}\left( \mathbf{x|y} \right)\), \emph{wobei{]}}

\[s_{i}(x) = y_{i} + b_{i} \cdot \left( x - x_{i} \right) + c_{i} \cdot \left( x - x_{i} \right)^{2} + d_{i} \cdot \left( x - x_{i} \right)^{3};i = 0,1,...,n - 1.\]

\begin{itemize}
\item
  Import of an \emph{ASCII} function \emph{matrix} file \(\mathbf{F}\);
\item
  Output of the \emph{ASCII} coefficient \emph{matrix} file
  \(\mathbf{A}\) (\textbf{\emph{KUSI.txt}});
\item
  Calculation of \(\mathbf{s}\left( \mathbf{x} \right)\) via the
  interpolation function.
\item
  \emph{{[}Übernahme einer ASCII Funktionsmatrix Datei}
  \(\mathbf{F}\)\emph{;{]}}
\item
  \emph{{[}Ausgabe der ASCII Koeffizientenmatrix Datei} \(\mathbf{A}\)
  \emph{(KUSI.txt);{]}}
\item
  \emph{{[}Berechnung von} \(\mathbf{s}\left( \mathbf{x} \right)\)
  \emph{über die Interpolations-Funktion.{]}}
\end{itemize}

Usage \emph{{[}Handhabung{]}}:

\begin{verbatim}
KUSI [f][x]
[f] ......... Function matrix file (F)
[x] ......... Function value x
\end{verbatim}

Example \emph{{[}Bsp.{]}}:

\begin{verbatim}
KUSI in.txt 0.5
\end{verbatim}

\hypertarget{kusf}{%
\subsubsection{3.4. KUSF}\label{kusf}}

\begin{center}\rule{0.5\linewidth}{0.5pt}\end{center}

\emph{Cubic spline} function: Calculation of a function \emph{matrix}
\(\mathbf{S =}\left( \mathbf{x|s}\left( \mathbf{x} \right) \right)\)
corresponding to coefficient \emph{matrix}
\(\mathbf{A =}\left( \mathbf{b}\left| \mathbf{c} \right|\mathbf{d} \right)\),
where \emph{{[}Kubische Spline Funktion: Berechnung einer
Funktionsmatrix\textless{}\ins\textgreater{}}
\(\mathbf{S =}\left( \mathbf{x|s}\left( \mathbf{x} \right) \right)\)
\emph{zu Koeffizientenmatrix}
\(\mathbf{A =}\left( \mathbf{b}\left| \mathbf{c} \right|\mathbf{d} \right)\),
\emph{wobei{]}}

\[s_{i}(x) = y_{i} + b_{i} \cdot \left( x - x_{i} \right) + c_{i} \cdot \left( x - x_{i} \right)^{2} + d_{i} \cdot \left( x - x_{i} \right)^{3};i = 0,1,...,n - 1.\]

\begin{itemize}
\item
  Import of an \emph{ASCII} function \emph{matrix} file \(\mathbf{F}\);
\item
  Import of the \emph{ASCII} coefficient \emph{matrix} file
  \(\mathbf{A}\) (\textbf{\emph{KUSI.txt}});
\item
  Output of the \emph{ASCII} function \emph{matrix} file \(\mathbf{S}\)
  (\textbf{\emph{KUSF.txt}}).
\item
  \emph{{[}Übernahme einer ASCII Funktionsmatrix Datei}
  \(\mathbf{F}\)\emph{;{]}}
\item
  \emph{{[}Übernahme der ASCII Koeffizientenmatrix Datei} \(\mathbf{A}\)
  \emph{(KUSI.txt);{]}}
\item
  \emph{{[}Ausgabe der ASCII Funktionsmatrix Datei} \(\mathbf{S}\)
  \emph{(KUSF.txt).{]}}
\end{itemize}

Usage \emph{{[}Handhabung{]}}:

\begin{verbatim}
KUSF [f][a][b][d] 
[f] ......... Function matrix file (F)
[a] ......... (x) Minimum
[b] ......... (x) Maximum
[d] ......... Interval d
\end{verbatim}

Example \emph{{[}Bsp.{]}}:

\begin{verbatim}
KUSF in.txt 1 10 0.1
\end{verbatim}

\hypertarget{nwti}{%
\subsubsection{3.5. NWTI}\label{nwti}}

\begin{center}\rule{0.5\linewidth}{0.5pt}\end{center}

\emph{Newton} interpolation: Calculation of the coefficient
\emph{vector} \(\mathbf{a}\) as well as
\(\mathbf{p}\left( \mathbf{x} \right)\) for an (empirical) function
\emph{matrix} \(\mathbf{F =}\left( \mathbf{x|y} \right)\), where
\emph{{[}Newton Interpolation: Berechnung des Koeffizientenvektors}
\(\mathbf{a}\) sowie \(\mathbf{p}\left( \mathbf{x} \right)\) \emph{zu
einer (empirischen) Funktionsmatrix}
\(\mathbf{F =}\left( \mathbf{x|y} \right)\), \emph{wobei{]}}

\[p(x) = a_{0} + a_{1} \cdot (x - 1) + a_{2} \cdot (x - 1) \cdot (x - 2)...a_{n} \cdot (x - 1) \cdot (x - 2) \cdot ... \cdot (x - n).\]

\begin{itemize}
\item
  Import of an \emph{ASCII} function \emph{matrix} file \(\mathbf{F}\);
\item
  Output of the \emph{ASCII} coefficient \emph{vector} file
  \(\mathbf{a}\) (\textbf{\emph{nwti.txt}});
\item
  Calculation of \(\mathbf{p}\left( \mathbf{x} \right)\) using the
  interpolation polynomial.
\item
  \emph{{[}Übernahme einer ASCII Funktionsmatrix Datei}
  \(\mathbf{F}\)\emph{;{]}}
\item
  \emph{{[}Ausgabe einer ASCII Koeffizientenvektor Datei} \(\mathbf{a}\)
  \emph{(nwti.txt);{]}}
\item
  \emph{{[}Berechnung von} \(\mathbf{p}\left( \mathbf{x} \right)\)
  \emph{über das Interpolations-Polynom.{]}}
\end{itemize}

Usage \emph{{[}Handhabung{]}}:

\begin{verbatim}
NWTI [f][x]
[f] ......... Function matrix file (F)
[x] ......... Function value x
\end{verbatim}

Example \emph{{[}Bsp.{]}}:

\begin{verbatim}
NWTI in.txt 3
\end{verbatim}

\hypertarget{nwtp}{%
\subsubsection{3.6. NWTP}\label{nwtp}}

\begin{center}\rule{0.5\linewidth}{0.5pt}\end{center}

\emph{Newton} Interpolation Polynomial: Calculation of a function
\emph{matrix}
\(\mathbf{F =}\left( \mathbf{x|p}\left( \mathbf{x} \right) \right)\)
corresponding to coefficient \emph{vector} \(\mathbf{a}\), where
\emph{{[}Newton Interpolations Polynom: Berechnung einer
Funktionsmatrix}
\(\mathbf{F =}\left( \mathbf{x|p}\left( \mathbf{x} \right) \right)\)
\emph{zu Koeffizientenvektor} \(\mathbf{a}\)\emph{, wobei{]}}

\[p(x) = a_{0} + a_{1} \cdot (x - 1) + a_{2} \cdot (x - 1) \cdot (x - 2)...a_{n} \cdot (x - 1) \cdot (x - 2) \cdot ... \cdot (x - n).\]

\begin{itemize}
\item
  Import of the \emph{ASCII} coefficient \emph{vector} file
  \(\mathbf{a}\) (\textbf{\emph{nwti.txt}});
\item
  Output of the \emph{ASCII} function \emph{matrix} file \(\mathbf{F}\)
  (\textbf{\emph{nwtp.txt}}).
\item
  \emph{{[}Übernahme der ASCII Koeffizientenvektor Datei} \(\mathbf{a}\)
  \emph{(nwti.txt);{]}}
\item
  \emph{{[}Ausgabe der ASCII Funktionsmatrix Datei} \(\mathbf{F}\)
  \emph{(nwtp.txt).{]}}
\end{itemize}

Usage \emph{{[}Handhabung{]}}:

\begin{verbatim}
NWTP [a][b][d]
[a] ......... (x) Minimum
[b] ......... (x) Maximum
[d] ......... Interval d
\end{verbatim}

Example \emph{{[}Bsp.{]}}:

\begin{verbatim}
NWTP 1 10 0.1
\end{verbatim}

\hypertarget{consoleapp_matrix}{%
\subsection{4. ConsoleApp\_Matrix}\label{consoleapp_matrix}}

Console applications for (4.1.-4.9.) \emph{matrix} calculation (Cayley,
\href{http://www.jstor.org/stable/108649}{1858}; c.f. Cardano,
\href{https://web.archive.org/web/20220201093634/http:/www.filosofia.unimi.it/cardano/testi/operaomnia/vol_4_s_4.pdf}{1545};
Heisenberg, \href{https://doi.org/10.1007/BF01328377}{1925}; Turing,
\href{https://turing.academicwebsite.com/publications/20-rounding-off-errors-in-matrix-processes}{1948}
or Crilly, \href{https://doi.org/10.1016/0315-0860(86)90091-1}{1986})
and (4.10.-4.21.) \emph{tools}, s. Schrausser
(\href{https://doi.org/10.5281/zenodo.7655046}{2023c}); on the topic in
general c.f. Meyberg and Vachenauer
(\href{https://doi.org/10.1007/978-3-642-56654-7}{2001}), Grcar
(\href{https://doi.org/10.1016/j.hm.2010.06.003}{2011}), Karpfinger
(\href{https://doi.org/10.1007/978-3-662-65458-3_10}{2022a},
\href{https://doi.org/10.1007/978-3-662-63305-2_11}{b}), Shores
(\href{https://doi.org/10.1007/978-3-319-74748-4_2}{2018}), Gentle
(\href{https://doi.org/10.1007/978-3-031-42144-0}{2024}) or Saff and
Snider (\href{https://doi.org/10.1007/978-3-031-97222-5}{2025}).

\hypertarget{ama}{%
\subsubsection{4.1. AMA}\label{ama}}

\begin{center}\rule{0.5\linewidth}{0.5pt}\end{center}

Adds or subtracts 2 \emph{matrices}, with given \emph{{[}Addiert oder
subtrahiert 2 Matrizen{]}},

\[k_{1} = k_{2}\text{,\ }n_{1} = n_{2}.\]

\emph{{[}wird vorausgesetzt{]}}.

\begin{itemize}
\item
  Import of 2 \emph{ASCII matrix} files.
\item
  Output of an \emph{ASCII matrix} file.
\item
  \emph{{[}Übernahme 2er ASCII Matrixdateien.{]}}
\item
  \emph{{[}Ausgabe einer ASCII Matrixdatei.{]}}
\end{itemize}

Usage \emph{{[}Handhabung{]}}:

\begin{verbatim}
AMA [matrix1][matrix2][output][mode] 
[matrix1] ... Input file 1 
[matrix2] ... Input file 2 
[output] .... Output file 
[mode] ...... (0):Addition (1):Subtraction
\end{verbatim}

Example \emph{{[}Bsp.{]}}:

\begin{verbatim}
AMA in1.txt in2.txt out.txt 0
\end{verbatim}

\hypertarget{ima}{%
\subsubsection{4.2. IMA}\label{ima}}

\begin{center}\rule{0.5\linewidth}{0.5pt}\end{center}

Calculates the inverse \(\mathbf{A}^{- 1}\) of \(\mathbf{A}\) using the
\emph{chained} form of \emph{Gaussian elimination {[}Berechnet die
inverse} \(\mathbf{A}^{- 1}\) \emph{von} \(\mathbf{A}\) \emph{über die
verkettete Form des Gaussschen Algorithmus{]}},

\[k_{\max} = n_{\max} = 200,a_{11} \neq 0.\]

This results in two triangular \emph{matrices} \(\mathbf{B}\) and
\(\mathbf{C}\), as well as the \emph{matrix} \(\mathbf{T}\) to the
generated identity \emph{matrix} \(\mathbf{E}\); \(\mathbf{A}^{- 1}\) is
obtained by its transpose \emph{{[}Es resultieren 2 Dreiecksmatrizen}
\(\mathbf{B}\) \emph{und} \(\mathbf{C}\)\emph{, sowie die Matrix}
\(\mathbf{T}\) \emph{zur erzeugten Einheitsmatrix}
\(\mathbf{E}\)\emph{,} \(\mathbf{A}^{- 1}\) \emph{entsteht
transponiert{]}:}

\begin{verbatim}
. . .  A   1 0 0  E
. . .      0 1 0
. . .      0 0 1

. . .  B   . . .  T
  . .      . . .
    .      . . .
.      C
. .

. . . (1/A)'
. . .
. . .
\end{verbatim}

\begin{itemize}
\item
  Import of a square \emph{ASCII matrix} file.
\item
  Output of a square \emph{ASCII matrix} file.
\item
  \emph{{[}Übernahme einer quadratischen ASCII Matrixdatei.{]}}
\item
  \emph{{[}Ausgabe einer quadratischen ASCII Matrixdatei.{]}}
\end{itemize}

Usage \emph{{[}Handhabung{]}}:

\begin{verbatim}
IMA [matrix][output] 
[matrix] ....... Input file 
[output] ....... Output file
\end{verbatim}

Example \emph{{[}Bsp.{]}}:

\begin{verbatim}
IMA in.txt out.txt
\end{verbatim}

\hypertarget{mma}{%
\subsubsection{4.3. MMA}\label{mma}}

\begin{center}\rule{0.5\linewidth}{0.5pt}\end{center}

Multiplies two \emph{matrices} (\(k_{1} = n_{2}\) is assumed). The
result is a \emph{matrix} with \emph{{[}Multipliziert 2 Matrizen
(}\(k_{1} = n_{2}\) \emph{wird vorausgesetzt). Es resultiert eine Matrix
mit{]}}

\[n = n_{1},\text{\ }k = k_{2}:\]

\begin{verbatim}
       *          =
 o . .   o o         o o
 o . .   . .         o .
         . .
 
 
 o . .   o o         o o
 o . .   . .         o .
 o . .   . .         o .    
 o . .               o .
 o . .               o .
\end{verbatim}

\begin{itemize}
\item
  Import of 2 \emph{ASCII} \emph{matrix} files.
\item
  Output of an \emph{ASCII matrix} file.
\item
  \emph{{[}Übernahme 2er ASCII Matrixdateien.{]}}
\item
  \emph{{[}Ausgabe einer ASCII Matrixdatei.{]}}
\end{itemize}

Usage \emph{{[}Handhabung{]}}:

\begin{verbatim}
MMA [matrix1][matrix2][output] 
[matrix1] ...... Input file 1 
[matrix2] ...... Input file 2 
[output] ....... Output file
\end{verbatim}

Example \emph{{[}Bsp.{]}}:

\begin{verbatim}
MMA m1.txt m2.txt m3.txt
\end{verbatim}

\hypertarget{qma}{%
\subsubsection{4.4. QMA}\label{qma}}

\begin{center}\rule{0.5\linewidth}{0.5pt}\end{center}

Squaring a square \emph{matrix {[}Quadriert eine quadratische
Matrix{]}}:

\begin{verbatim}
        *       =
 o .   o o      o o
 o .   . .      o .
\end{verbatim}

\begin{itemize}
\item
  Import of a square \emph{ASCII matrix} file.
\item
  Output of a square \emph{ASCII matrix} file.
\item
  \emph{{[}Übernahme einer quadratischen ASCII Matrixdateien.{]}}
\item
  \emph{{[}Ausgabe einer quadratischen ASCII Matrixdatei.{]}}
\end{itemize}

Usage \emph{{[}Handhabung{]}}:

\begin{verbatim}
QMA [matrix][output] 
[matrix] ....... Input file 
[output] ....... Output file
\end{verbatim}

Example \emph{{[}Bsp.{]}}:

\begin{verbatim}
QMA a.txt b.txt
\end{verbatim}

\hypertarget{sma}{%
\subsubsection{4.5. SMA}\label{sma}}

\begin{center}\rule{0.5\linewidth}{0.5pt}\end{center}

Selects a \emph{submatrix} (or \emph{vector}) from a \emph{matrix}
\emph{{[}Selegiert eine Sub Matrix (oder einen Vektor) aus einer
Matrix{]}}.

\begin{itemize}
\item
  Importing an \emph{ASCII matrix} file.
\item
  Output of an \emph{ASCII} (\emph{matrix}) file.
\item
  \emph{{[}Übernahme einer ASCII Matrixdatei.{]}}
\item
  \emph{{[}Ausgabe einer ASCII (Matrix-)Datei.{]}}
\end{itemize}

Usage \emph{{[}Handhabung{]}}:

\begin{verbatim}
SMA [matrix][output][i0][i1][j0][j1] 
[matrix] ...... Matrix file 
[output] ...... Matrix output file 
[i0] .......... From line 
[i1] .......... To line 
[j0] .......... From column 
[j1] .......... To column
\end{verbatim}

Example \emph{{[}Bsp.{]}}:

\begin{verbatim}
SMA in.txt out1.txt 1 5 2 2 //Columnvector j2
SMA in.txt out2.txt 2 2 1 5 //Linevector i2
SMA in.txt out3.txt 2 4 2 4 //Submatrix
SMA in.txt out4.txt 3 3 3 3 //Single number
\end{verbatim}

\hypertarget{spur}{%
\subsubsection{4.6. SPUR}\label{spur}}

\begin{center}\rule{0.5\linewidth}{0.5pt}\end{center}

Calculates the trace (\emph{sp}) of a square \emph{matrix}
\(\mathbf{A}\) \emph{{[}Berechnet die Spur (sp) einer quadratischen
Matrix} \(\mathbf{A}\)\emph{{]}}:

\begin{verbatim}
o . .  A
. o .  
. . o
       sp A
\end{verbatim}

\begin{itemize}
\item
  Import of a square \emph{ASCII matrix} file.
\item
  Output of \emph{sp} \(\mathbf{A}\) in file \textbf{\emph{SPUR.txt}}.
\item
  \emph{{[}Übernahme einer quadratischen ASCII Matrixdatei.{]}}
\item
  \emph{{[}Ausgabe von sp} \(\mathbf{A}\) \emph{in die Datei
  SPUR.txt.{]}}
\end{itemize}

Usage \emph{{[}Handhabung{]}}:

\begin{verbatim}
SPUR [matrix][mode] 
[matrix] ...... Input file 
[mode] ........ Mode of trace calculation: 
                (0):Addition (Standard) 
                (1):Multiplication (s. VMA.exe)
                (2):Subtraction
                (3):Division
\end{verbatim}

Example \emph{{[}Bsp.{]}}:

\begin{verbatim}
SPUR in.txt 0
\end{verbatim}

\hypertarget{trp}{%
\subsubsection{4.7. TRP}\label{trp}}

\begin{center}\rule{0.5\linewidth}{0.5pt}\end{center}

Transposes a data \emph{matrix} \emph{{[}Transponiert eine
Datenmatrix{]}},

\[n_{\max} = k_{\max} = 1299.\]

Column separator, input file: tab or space. Column separator, output
file: 1 space \emph{{[}Spaltentrennzeichen, Eingabedatei: Tabulator oder
Leerzeichen. Spaltentrennzeichen, Ausgabedatei: 1 Leerzeichen{]}.}

\begin{itemize}
\item
  Importing an \emph{ASCII} data \emph{matrix} file.
\item
  Output of a transposed \emph{ASCII} data \emph{matrix} file.
\item
  \emph{{[}Übernahme einer ASCII Datenmatrixdatei.{]}}
\item
  \emph{{[}Ausgabe einer transponierten ASCII Datenmatrixdatei.{]}}
\end{itemize}

Usage \emph{{[}Handhabung{]}}:

\begin{verbatim}
TRP [input][output] 
[input] ....... Input file 
[output] ...... Output file
\end{verbatim}

Example \emph{{[}Bsp.{]}}:

\begin{verbatim}
TRP in.txt out.txt
\end{verbatim}

\hypertarget{vma}{%
\subsubsection{4.8. VMA}\label{vma}}

\begin{center}\rule{0.5\linewidth}{0.5pt}\end{center}

Calculates the \emph{chained} form of \emph{Gaussian elimination} of a
square \emph{matrix} \(\mathbf{A}\), with \emph{{[}Berechnet die
verkettete Form des Gaussschen Algorithmus einer quadratischen Matrix}
\(\mathbf{A}\)\emph{, mit{]}}

\[k_{\max} = n_{\max} = 250,a_{11} \neq 0.\]

This results in two triangular matrices \(\mathbf{B}\) and
\(\mathbf{C}\) \emph{{[}Es resultieren 2 Dreiecksmatrizen}
\(\mathbf{B}\) \emph{und} \(\mathbf{C}\)\emph{{]}}:

\begin{verbatim}
. . .  A
. . .  
. . .

. . .  B
  . .
    .
.      C
. .
\end{verbatim}

The determinant of \(\mathbf{A}\) (det \(\mathbf{A}\)) is the product of
the elements in the main diagonal of \(\mathbf{B}\)
(\(\mathbf{TT}b_{ii}\)) \emph{{[}Die Determinante von} \(\mathbf{A}\)
\emph{(det} \(\mathbf{A})\) \emph{ist das Produkt der Elemente in der
Hauptdiagonale von} \(\mathbf{B}\)
\emph{(}\(\mathbf{TT}b_{ii}\)\emph{){]}}.

\begin{itemize}
\item
  Import of a square \emph{ASCII matrix} file.
\item
  Output of a square \emph{ASCII matrix} file.
\item
  \emph{{[}Übernahme einer quadratischen ASCII Matrixdatei.{]}}
\item
  \emph{{[}Ausgabe einer quadratischen ASCII Matrixdatei.{]}}
\end{itemize}

Usage \emph{{[}Handhabung{]}}:

\begin{verbatim}
VMA [matrix][output] 
[matrix] ...... Input file 
[output] ...... Output file
\end{verbatim}

Example \emph{{[}Bsp.{]}}:

\begin{verbatim}
VMA in.txt out.txt
\end{verbatim}

\hypertarget{zma}{%
\subsubsection{4.9. ZMA}\label{zma}}

\begin{center}\rule{0.5\linewidth}{0.5pt}\end{center}

Multiplies a \emph{matrix} by a real number \emph{{[}Multipliziert eine
Matrix mit einer reellen Zahl{]}}.

\begin{itemize}
\item
  Import of an \emph{ASCII matrix} file.
\item
  Output of an \emph{ASCII matrix} file.
\item
  \emph{{[}Übernahme einer ASCII Matrixdatei.{]}}
\item
  \emph{{[}Ausgabe einer ASCII Matrixdatei.{]}}
\end{itemize}

Usage \emph{{[}Handhabung{]}}:

\begin{verbatim}
ZMA [matrix][output][wert] 
[matrix] ...... Input file 
[output] ...... Output file 
[wert] ........ Real number
\end{verbatim}

Example \emph{{[}Bsp.{]}}:

\begin{verbatim}
ZMA in.txt out.txt 3
\end{verbatim}

\hypertarget{ent}{%
\subsubsection{4.10. ENT}\label{ent}}

\begin{center}\rule{0.5\linewidth}{0.5pt}\end{center}

Performs a symmetrical interlaced split of a data \emph{vector} file
\(\mathbf{x_0}\) \emph{{[}Führt eine symmetrische entwobene Aufteilung
einer Datenvektordatei} \(\mathbf{x_0}\) \emph{durch{]}}:

\begin{verbatim}
x0
--
1
2
3
4

x1 x2
-- --
1
   2
3
   4
\end{verbatim}

\begin{itemize}
\item
  Importing a single-column \emph{ASCII} file.
\item
  Output of 2 single-column \emph{ASCII} files.
\item
  \emph{{[}Übernahme einer einspaltigen ASCII Datei.{]}}
\item
  \emph{{[}Ausgabe von 2 einspaltigen ASCII Dateien.{]}}
\end{itemize}

Usage \emph{{[}Handhabung{]}}:

\begin{verbatim}
ENT [input][output1][output2] 
[input] ...... Input file  
[output1] .... Output file 1 
[output2] .... Output file 2
\end{verbatim}

Example \emph{{[}Bsp.{]}}:

\begin{verbatim}
ENT in.txt out1.txt out2.txt
\end{verbatim}

\hypertarget{ktf}{%
\subsubsection{4.11. KTF}\label{ktf}}

\begin{center}\rule{0.5\linewidth}{0.5pt}\end{center}

Reduces or increases the size of a \emph{perfectly} linear data
\emph{vector}. The data adjustment up to \(n'\) is iteratively performed
by \emph{{[}Verringert oder vergrössert den Umfang eines perfekt
linearen Datenvektors. Die bis} \(n'\) \emph{iterative Datenanpassung
erfolgt über{]}}

\[x_{i}\lbrack n\rbrack = x_{i}\lbrack n + 1\rbrack \cdot \frac{n}{n - 1};n' < n,\]
\[x_{i}\lbrack n\rbrack = x_{i}\lbrack n - 1\rbrack \cdot \frac{n - 2}{n - 1};n' > n.\]

\begin{itemize}
\item
  Import of a single-column, ascendingly ordered \emph{ASCII} data
  \emph{vector} file of size \(n\).
\item
  Output of a single-column, ascendingly ordered \emph{ASCII} data
  \emph{vector} file of size \(n'\).
\item
  \emph{{[}Übernahme einer einspaltigen, aufsteigend geordneten ASCII
  Datenvektordatei im Umfang} \(n.\)\emph{{]}}
\item
  \emph{{[}Ausgabe einer einspaltigen, aufsteigend geordneten ASCII
  Datenvektordatei im Umfang} \(n^{'}.\)\emph{{]}}
\end{itemize}

Usage \emph{{[}Handhabung{]}}:

\begin{verbatim}
ktf [input][output][n] 
[input] ....... Input file  
[output] ...... Output file 
[n] ........... Vector size n’
\end{verbatim}

Example \emph{{[}Bsp.{]}}:

\begin{verbatim}
ktf data_in.txt data_out.txt 30
\end{verbatim}

\hypertarget{ktf2}{%
\subsubsection{4.12. KTF2}\label{ktf2}}

\begin{center}\rule{0.5\linewidth}{0.5pt}\end{center}

Reduces or increases the size of a data \emph{vector}
\emph{{[}Verringert oder vergrössert den Umfang eines Datenvektors{]}}
(\(n_{\max} = n'_{\max} = 33000\)).

\begin{itemize}
\item
  Import of a single-column, ascendingly ordered \emph{ASCII} data
  \emph{vector} file of size \(n\).
\item
  Output of a single-column, ascendingly ordered \emph{ASCII} data
  \emph{vector} file of size \(n'\).
\item
  \emph{{[}Übernahme einer einspaltigen, aufsteigend geordneten ASCII
  Datenvektordatei im Umfang} \(n.\)\emph{{]}}
\item
  \emph{{[}Ausgabe einer einspaltigen, aufsteigend geordneten ASCII
  Datenvektordatei im Umfang} \(n^{'}.\)\emph{{]}}
\end{itemize}

Usage \emph{{[}Handhabung{]}}:

\begin{verbatim}
ktf2 [input][output][n] 
[input] ....... Input file 
[output] ...... Output file 
[n] ........... Vector size n’
\end{verbatim}

Example \emph{{[}Bsp.{]}}:

\begin{verbatim}
ktf2 in.txt out.txt 12
ktf2 in.txt out.txt 110
\end{verbatim}

\hypertarget{ktf3}{%
\subsubsection{4.13. KTF3}\label{ktf3}}

\begin{center}\rule{0.5\linewidth}{0.5pt}\end{center}

Adapts a data \emph{vector} to a target coordinate system. The data
adaptation is performed by \emph{{[}Passt einen Datenvektor an ein
Ziel-Koordinatensystem an. Die Datenanpassung erfolgt über{]}}

\[x_{i}' = min_{x} + \{\lbrack\left( min_{x} - x_{\min} \right) - \left( min_{x} - x_{i} \right)\rbrack \cdot \frac{max_{x} - min_{x}}{\left( min_{x} - x_{\min} \right) - \left( min_{x} - x_{\max} \right)}\}\]

When inverting a value, \(x''_i\) is calculated as \emph{{[}bei einer
Wertinvertierung errechnet man} \(x''_i\) \emph{über{]}}

\[x''_{i} = \left( min_{x} + max_{x} \right) - x'_{i},\]

where \emph{{[}mit{]}}

\(min_{x}\) \ldots.. Value of the minimum point in the target coordinate
system\\
\(max_{x}\) \ldots. Value of the maximum point in the target coordinate
system\\
\(x_{\min}\) \ldots\ldots. \emph{Vector} minimum value\\
\(x_{\max}\) \ldots... \emph{Vector} maximum value

\(min_{x}\) \ldots.. \emph{{[}Wert des Minimalpunktes im
Ziel-Koordinatensystem{]}}\\
\(max_{x}\) \ldots.\emph{{[}Wert des Maximalpunktes im
Ziel-Koordinatensystem{]}}\\
\(x_{\min}\) \ldots\ldots. \emph{{[}Vektor Minimalwert{]}}\\
\(x_{\max}\) \ldots\ldots{} \emph{{[}Vektor Maximalwert{]}}

\begin{itemize}
\item
  Import of a single-column, ascendingly ordered \emph{ASCII} data
  \emph{vector} file.
\item
  Output of a two-column, ascendingly ordered \emph{ASCII} data
  \emph{matrix} file containing:
\item
  \emph{{[}Übernahme einer einspaltigen, aufsteigend geordneten ASCII
  Datenvektordatei.{]}}
\item
  \emph{{[}Ausgabe einer zweispaltigen, aufsteigend geordneten ASCII
  Datenmatrixdatei beinhaltend:{]}}
\end{itemize}

The data \emph{vector} adapted to the target coordinate system; the
initial data \emph{vector} \emph{{[}Den an das Ziel-Koordinatensystem
angepassten Datenvektor; den ursprünglichen Datenvektor{]}}.

Usage \emph{{[}Handhabung{]}}:

\begin{verbatim}
ktf3 [input][output][minx][maxx][inv] 
[input] ....... Input file 
[output] ...... Output file 
[minx] ........ Val. of the min. point in the target coordinate system
[maxx] ........ Val. of the max. point in the target coordinate system
[inv] ......... (1):Value inversion (0):No Value inversion
\end{verbatim}

Example \emph{{[}Bsp.{]}}:

\begin{verbatim}
ktf3 data_in.txt data_out.txt 1 100 1
ktf3 data_in.txt data_out.txt 371 51 0
\end{verbatim}

\hypertarget{ntf}{%
\subsubsection{4.14. NTF}\label{ntf}}

\begin{center}\rule{0.5\linewidth}{0.5pt}\end{center}

Generates an ascending linear data \emph{vector} and fits it to a target
coordinate system. The data fitting is performed by \emph{{[}Erzeugt
einen aufsteigend geordneten linearen Datenvektor und passt diesen an
ein Ziel-Koordinatensystem an. Die Datenanpassung erfolgt über{]}}

\[x_{i}' = min_{x} + \{\lbrack\left( min_{x} - x_{\min} \right) - \left( min_{x} - x_{i} \right)\rbrack \cdot \frac{max_{x} - min_{x}}{\left( min_{x} - x_{\min} \right) - \left( min_{x} - x_{\max} \right)}\}\]

When inverting a value, \(x_{i}''\) is calculated as \emph{{[}bei einer
Wertinvertierung errechnet man} \(x_{i}''\) \emph{über{]}}

\[x''_{i} = \left( min_{x} + max_{x} \right) - x'_{i},\]

where \emph{{[}mit{]}}

\(min_{x}\) \ldots.. Value of the minimum point in the target coordinate
system\\
\(max_{x}\) \ldots. Value of the maximum point in the target coordinate
system\\
\(x_{\min}\) \ldots\ldots. \emph{Vector} minimum value\\
\(x_{\max}\) \ldots\ldots{} \emph{Vector} maximum value

\(min_{x}\) \ldots.. \emph{{[}Wert des Minimalpunktes im
Ziel-Koordinatensystem{]}}\\
\(max_{x}\) \ldots. \emph{{[}Wert des Maximalpunktes im
Ziel-Koordinatensystem{]}}\\
\(x_{\min}\) \ldots\ldots. \emph{{[}Vektor Minimalwert{]}}\\
\(x_{\max}\) \ldots\ldots{} \emph{{[}Vektor Maximalwert{]}}

\begin{itemize}
\tightlist
\item
  Output of a single-column, ascending linear \emph{ASCII} data
  \emph{vector} file.
\item
  \emph{{[}Ausgabe einer einspaltigen, aufsteigend geordneten linearen
  ASCII Datenvektordatei.{]}}
\end{itemize}

Usage \emph{{[}Handhabung{]}}:

\begin{verbatim}
ntf [output][minn][maxn][min][max][inv] 
[output] ...... Output file
[minn] ........ n Minimum value 
[maxn] ........ n Maximum value 
[min] ......... Coordinates Minimum Position Value 
[max] ......... Coordinates Maximum Position Value 
[inv] ......... (1):inverted (0):not inverted
\end{verbatim}

Example \emph{{[}Bsp.{]}}:

\begin{verbatim}
ntf data_out.txt 1 30000 1 30000 0
ntf data_out.txt 1 30000 372 50 1
\end{verbatim}

\hypertarget{sel}{%
\subsubsection{4.15. SEL}\label{sel}}

\begin{center}\rule{0.5\linewidth}{0.5pt}\end{center}

Selects a data \emph{vector} from a data \emph{matrix}
\emph{{[}Selegiert einen Datenvektor aus einer Datenmatrix{]}},

\[n_{\max} = 33000.\]

\begin{itemize}
\item
  Import of an \emph{ASCII} data \emph{matrix} file.
\item
  Output of a single-column \emph{ASCII} data \emph{vector} file.
\item
  \emph{{[}Übernahme einer ASCII Datenmatrixdatei.{]}}
\item
  \emph{{[}Ausgabe einer einspaltigen ASCII Datenvektordatei.{]}}
\end{itemize}

Usage \emph{{[}Handhabung{]}}:

\begin{verbatim}
sel [input][output][a][k] 
[input] ....... Input file 
[output] ...... Output file 
[a] ........... Vector number 
[k] ........... Number of vectors
\end{verbatim}

Example \emph{{[}Bsp.{]}}:

\begin{verbatim}
sel in.txt out.asc 3 5
\end{verbatim}

\hypertarget{srt}{%
\subsubsection{4.16. SRT}\label{srt}}

\begin{center}\rule{0.5\linewidth}{0.5pt}\end{center}

Sorts a data \emph{vector} \emph{{[}Sortiert einen Datenvektor{]}},

\[n_{\max} = 33000,\]

16-digit output \emph{{[}16 stellige Ausgabe{]}}.

\begin{itemize}
\item
  Import of a single-column \emph{ASCII} data \emph{vector} file.
\item
  Output of a sorted single-column \emph{ASCII} data \emph{vector} file.
\item
  \emph{{[}Übernahme einer einspaltigen ASCII Datenvektordatei.{]}}
\item
  \emph{{[}Ausgabe einer sortierten einspaltigen ASCII
  Datenvektordatei.{]}}
\end{itemize}

Usage \emph{{[}Handhabung{]}}:

\begin{verbatim}
srt [input][output] [[d]] 
[input] ...... Input file 
[output] ..... Output file 
[d] .......... optional (1):descending sorting
\end{verbatim}

Example \emph{{[}Bsp.{]}}:

\begin{verbatim}
srt in.txt out.txt
srt in.txt out.txt 1
\end{verbatim}

\hypertarget{srt1}{%
\subsubsection{4.17. SRT1}\label{srt1}}

\begin{center}\rule{0.5\linewidth}{0.5pt}\end{center}

Concatenates two sorted data \emph{vectors} using filestream processing
\emph{{[}Verkettet 2 sortierte Datenvektoren mit Filestream
Verarbeitung{]}},

\[n_{\max} \rightarrow \infty.\]

\begin{itemize}
\item
  Import of two ascendingly sorted single-column \emph{ASCII} data
  \emph{vector} files.
\item
  Output of a sorted single-column \emph{ASCII} file \emph{vector} file.
\item
  \emph{{[}Übernahme von zwei aufsteigend sortierten einspaltigen ASCII
  Datenvektordateien.{]}}
\item
  \emph{{[}Ausgabe einer sortierten einspaltigen ASCII
  Datenvektordatei.{]}}
\end{itemize}

Usage \emph{{[}Handhabung{]}}:

\begin{verbatim}
srt1 [input1][input2][output] 
[input1] ...... Input file 1 
[input2] ...... Input file 2 
[output] ...... Output file
\end{verbatim}

Example \emph{{[}Bsp.{]}}:

\begin{verbatim}
srt1 in1.asc in2.asc out.asc
\end{verbatim}

\hypertarget{srt2}{%
\subsubsection{4.18. SRT2}\label{srt2}}

\begin{center}\rule{0.5\linewidth}{0.5pt}\end{center}

Sorts a data \emph{vector {[}Sortiert einen Datenvektor{]}},

\[n_{\max} = 33000.\]

Data sorting is performed via iterative pairwise comparison \emph{{[}Die
Datensortierung erfolgt über iterativen Paarvergleich{]}}

\[i\text{\ vs.\ }i + 1\]

and pairwise exchange \emph{{[}und Paartausch{]}}

\[i > i + 1.\]

(slower than \textbf{\emph{SRT.EXE}}) \emph{{[}(langsamer als
SRT.EXE){]}}.

\begin{itemize}
\item
  Import of a single-column \emph{ASCII} data \emph{vector} file.
\item
  Output of a sorted single-column \emph{ASCII} data \emph{vector} file.
\item
  \emph{{[}Übernahme einer einspaltigen ASCII Datenvektordatei.{]}}
\item
  \emph{{[}Ausgabe einer sortierten einspaltigen ASCII
  Datenvektordatei.{]}}
\end{itemize}

Usage \emph{{[}Handhabung{]}}:

\begin{verbatim}
srt2 [input][output] [[d]] 
[input] ...... Input file 
[output] ..... Output file
[d] .......... optional (1):descending sorting
\end{verbatim}

Example \emph{{[}Bsp.{]}}:

\begin{verbatim}
srt2 in.txt out.txt
srt2 in.txt out.txt 1
\end{verbatim}

\hypertarget{srt3}{%
\subsubsection{4.19. SRT3}\label{srt3}}

\begin{center}\rule{0.5\linewidth}{0.5pt}\end{center}

Sorts a data \emph{vector} \emph{{[}Sortiert einen Datenvektor{]}},

\[n_{\max} = 33000,\]

max. 8 digits \emph{{[}max. 8-Stellen{]}}.

Very fast calculation through implementation of the \emph{C}-function
\textbf{\emph{Qsort}} \emph{{[}Sehr schnelle Berechnung durch Umsetzung
der C-eigenen Qsort Funktion{]}}.

\begin{itemize}
\item
  Import of a single-column \emph{ASCII} data \emph{vector} file.
\item
  Output of a sorted single-column \emph{ASCII} data \emph{vector} file.
\item
  \emph{{[}Übernahme einer einspaltigen ASCII Datenvektordatei.{]}}
\item
  \emph{{[}Ausgabe einer sortierten einspaltigen ASCII
  Datenvektordatei.{]}}
\end{itemize}

Usage \emph{{[}Handhabung{]}}:

\begin{verbatim}
srt3 [input][output] [[d]] 
[input] ...... Input file 
[output] ..... Output file 
[d] .......... optional (1):descending sorting
\end{verbatim}

Example \emph{{[}Bsp.{]}}:

\begin{verbatim}
srt3 in.txt out.txt
srt3 in.txt out.txt 1
\end{verbatim}

\hypertarget{v2v}{%
\subsubsection{4.20. V2V}\label{v2v}}

\begin{center}\rule{0.5\linewidth}{0.5pt}\end{center}

Concatenates two single-column \emph{ASCII} files (\(n_{1} = n_{2}\)
required) \emph{{[}Fügt 2 einspaltige ASCII Dateien aneinander
(}\(n_{1} = n_{2}\) \emph{wird vorausgesetzt){]}}.

\begin{itemize}
\item
  Import of 2 single-column \emph{ASCII} files.
\item
  Output of a two-column \emph{ASCII} file.
\item
  \emph{{[}Übernahme von 2 einspaltigen ASCII Dateien.{]}}
\item
  \emph{{[}Ausgabe einer zweispaltigen ASCII Datei.{]}}
\end{itemize}

Usage \emph{{[}Handhabung{]}}:

\begin{verbatim}
V2V [input1][input2][output][tab] 
[input1] ...... Input file 1 
[input2] ...... Input file 2 
[output] ...... Output file 
[tab] ......... Column separator
                (0):Tabulator 
                (1):Blank 
                (*):String
\end{verbatim}

Example \emph{{[}Bsp.{]}}:

\begin{verbatim}
V2V in1.txt in2.txt out.txt 1
V2V in1.txt in2.txt out.txt <---->
\end{verbatim}

\hypertarget{z2z}{%
\subsubsection{4.21. Z2Z}\label{z2z}}

\begin{center}\rule{0.5\linewidth}{0.5pt}\end{center}

Concatenates two \emph{ASCII} files \emph{{[}Fügt zwei ASCII Dateien
aneinander{]}}.

\begin{itemize}
\item
  Import of two \emph{ASCII} files.
\item
  Output of an \emph{ASCII} file.
\item
  \emph{{[}Übernahme zweier ASCII Dateien.{]}}
\item
  \emph{{[}Ausgabe einer ASCII Datei.{]}}
\end{itemize}

Usage \emph{{[}Handhabung{]}}:

\begin{verbatim}
Z2Z [input1][input2][output] 
[input1] ...... Input file 1 
[input2] ...... Input file 2 
[output] ...... Output file 
\end{verbatim}

Example \emph{{[}Bsp.{]}}:

\begin{verbatim}
Z2Z in1.txt in2.txt out.txt
\end{verbatim}

\textbf{Table 2}. \emph{ConsoleApp} functions overview with
corresponding functions in \texttt{SCHRAUSSER-MAT} (Schrausser,
\href{https://doi.org/10.17605/OSF.IO/8XE42}{2022}).

\begin{verbatim}
Category                    chpt.   nr.     Function    
                                            ConsoleApp      SCHRAUSSER-MAT

EEG/ Focus parameters       2       1       ERC, ERCX   
                            2       2       XCOR    
                            2       3       FOC¹    
                            2       4       DIS¹    
                            2       5       FLOC¹   
                            2       6       OUT¹    

Integral/ Interpolation     3       1       ROMI    
                            3       2       ROME    
                            3       3       KUSI    
                            3       4       KUSF    
                            3       5       NWTI    
                            3       6       NWTP    

Matrix/ Data manipulation   4       1       AMA             MXA, MXS
                            4       2       IMA             MXI
                            4       3       MMA             MXM
                            4       4       QMA             MXQ
                            4       5       SMA 
                            4       6       SPUR            MSP
                            4       7       TRP             MXT
                            4       8       VMA             MXV
                            4       9       ZMA 
                            4       10      ENT 
                            4       11      KTF 
                            4       12      KTF2    
                            4       13      KTF3    
                            4       14      NTF 
                            4       15      SEL 
                            4       16      SRT 
                            4       17      SRT1    
                            4       18      SRT2    
                            4       19      SRT3    
                            4       20      V2V 
                            4       21      Z2Z      
\end{verbatim}

¹) implemented in \texttt{ATINA} (Schrausser,
\href{https://doi.org/10.5281/zenodo.19422127}{2006}).

\hypertarget{consoleapp_string}{%
\subsection{5. ConsoleApp\_String}\label{consoleapp_string}}

Console applications for \emph{string-} and
\emph{number-transformation}, see Schrausser
(\href{https://doi.org/10.5281/zenodo.7653790}{2023d}).

\hypertarget{dec}{%
\subsubsection{5.1. 2DEC}\label{dec}}

\begin{center}\rule{0.5\linewidth}{0.5pt}\end{center}

Converts the argument to a decimal \emph{number} \emph{{[}Argument in
Dezimalzahl{]}}.

Usage \emph{{[}Handhabung{]}}:

\begin{verbatim}
2DEC [x] 
[x] .......... Argument
\end{verbatim}

\hypertarget{dec2}{%
\subsubsection{5.2. DEC2}\label{dec2}}

\begin{center}\rule{0.5\linewidth}{0.5pt}\end{center}

Decimal \emph{number} conversion \emph{{[}Dezimalzahl Umwandlung{]}}.

Usage \emph{{[}Handhabung{]}}:

\begin{verbatim}
DEC2 [d]
[d] .......... Decimal number
\end{verbatim}

\hypertarget{bin2dec}{%
\subsubsection{5.3. BIN2DEC}\label{bin2dec}}

\begin{center}\rule{0.5\linewidth}{0.5pt}\end{center}

Binary \emph{number} converted to decimal \emph{number}
\emph{{[}Binärzahl in Dezimalzahl{]}}.

Usage \emph{{[}Handhabung{]}}:

\begin{verbatim}
BIN2DEC [b]
[b] .......... Binary number
\end{verbatim}

\hypertarget{dec2bin}{%
\subsubsection{5.4. DEC2BIN}\label{dec2bin}}

\begin{center}\rule{0.5\linewidth}{0.5pt}\end{center}

Binary \emph{number} converted to decimal \emph{number}
\emph{{[}Dezimalzahl in Binärzahl{]}}.

Usage \emph{{[}Handhabung{]}}:

\begin{verbatim}
DEC2BIN [d]
[d] .......... Decimal number
\end{verbatim}

\hypertarget{abcd}{%
\subsubsection{5.5. ABCD}\label{abcd}}

\begin{center}\rule{0.5\linewidth}{0.5pt}\end{center}

Modifies or deletes \emph{strings} in an \emph{ASCII} file
\emph{{[}Ändert oder löscht Strings einer ASCII Datei{]}.}

\begin{itemize}
\item
  Import of an \emph{ASCII} file.
\item
  Output of a modified \emph{ASCII} file.
\item
  \emph{{[}Übernahme einer ASCII Datei.{]}}
\item
  \emph{{[}Ausgabe einer modifizierten ASCII Datei.{]}}
\end{itemize}

Usage \emph{{[}Handhabung{]}}:

\begin{verbatim}
ABCD [input][output][ab][cd][sw]
[input] ...... Input file
[output] ..... Output file 
[ab] ......... String to be modified
[cd] ......... String to be inserted
[sw] ......... (0):modifies 
               (1):deletes (2):inserts a space
\end{verbatim}

Example \emph{{[}Bsp.{]}}:

\begin{verbatim}
ABCD input.txt output.txt .txt .asc 0
ABCD input.txt output.txt REM 0 1
\end{verbatim}

\hypertarget{azube}{%
\subsubsection{5.6. AZUBE}\label{azube}}

\begin{center}\rule{0.5\linewidth}{0.5pt}\end{center}

Changing individual characters in an \emph{ASCII} file \emph{{[}Änderung
einzelner Zeichen einer ASCII Datei{]}.}

Usage \emph{{[}Handhabung{]}}:

\begin{verbatim}
AZUBE [input][output][a][b] 
[input] ...... Input file
[output] ..... Output file
[a] .......... Character
[b] .......... Modify 
\end{verbatim}

Example \emph{{[}Bsp.{]}}:

\begin{verbatim}
AZUBE input.txt output.txt , .
\end{verbatim}

\hypertarget{tezute}{%
\subsubsection{5.7. TEZUTE}\label{tezute}}

\begin{center}\rule{0.5\linewidth}{0.5pt}\end{center}

Changing individual \emph{strings} in a file \emph{{[}Änderung einzelner
Strings einer Datei{]}.}

Usage \emph{{[}Handhabung{]}}:

\begin{verbatim}
TEZUTE [input][output][a][b] 
[input] ...... Input file
[output] ..... Output file 
[a] .......... String
[b] .......... Modify 
\end{verbatim}

Example \emph{{[}Bsp.{]}}:

\begin{verbatim}
78 00
67 8
null 00
9 null
\end{verbatim}

\begin{verbatim}
TEZUTE input.txt output.txt null 00
\end{verbatim}

\begin{verbatim}
78
00
67
8
00
00
9
00
\end{verbatim}

Edit source code for formatted output \emph{{[}Zur formatierten Ausgabe
Quellcode bearbeiten{]}}.

\hypertarget{ab_l}{%
\subsubsection{5.8. AB\_L}\label{ab_l}}

\begin{center}\rule{0.5\linewidth}{0.5pt}\end{center}

Inserts a \emph{string} (i) before and/or a \emph{string} (ii) after
each line of a file \emph{{[}Fügt einen String (i) vor und/oder einen
String (ii) nach jeder Zeile einer Datei ein{]}}.

\begin{itemize}
\item
  Import of a single-column \emph{ASCII} file.
\item
  Output of a (one-,) two- or three-column \emph{ASCII} file.
\item
  \emph{{[}Übernahme einer einspaltigen ASCII Datei.{]}}
\item
  \emph{{[}Ausgabe einer (ein-,) zwei- oder dreispaltigen ASCII
  Datei.{]}}
\end{itemize}

Usage \emph{{[}Handhabung{]}}:

\begin{verbatim}
AB_L [input][output][a][sw][z][sw]
[input] ...... Input file 
[output] ..... Output file
[a] .......... String inserted before each line of input file
[sw] ......... (0):no insert before each line 
[z] .......... String inserted after each line of input file
[sw] ......... (0):no insert after each line
\end{verbatim}

Example \emph{{[}Bsp.{]}}:

\begin{verbatim}
AB_L in.txt out.txt move 1 verzeichnis 1
AB_L in.txt out.txt del 1 0 0
\end{verbatim}

\hypertarget{biia}{%
\subsubsection{5.9. BIIA}\label{biia}}

\begin{center}\rule{0.5\linewidth}{0.5pt}\end{center}

Encodes a file into an \emph{ASCII} file (recoding via
\textbf{\emph{AIIB.EXE}}) \emph{{[}Encodiert eine Datei in eine ASCII
Datei (Recodierung via AIIB.EXE){]}}.

\begin{itemize}
\item
  Import of a file.
\item
  Output of an encoded \emph{ASCII} file (\textbf{\emph{*.bii}}).
\item
  \emph{{[}Übernahme einer Datei.{]}}
\item
  \emph{{[}Ausgabe einer} \emph{encodierten ASCII Datei (*.bii).{]}}
\end{itemize}

Usage \emph{{[}Handhabung{]}}:

\begin{verbatim}
BIIA [input][output][mode][form]
[input] ...... Input file
[output] ..... Output file in ASCII format (*.bii)
[mode] ....... (1):cd01 SCHRAUSSER CODE, uppercase letter combination (2 characters, e.g. AA, FD, ..., good readability – unambiguous, for archiving)
               (2):cd02 SCHRAUSSER CODE, keyboard character comb. (1 or 2 characters, e.g. a, 1A, ..., small file) 
               (3):cd03 SCHRAUSSER CODE, number combination (3 characters, e.g. 001, 123, ..., optimal readability)
               (4):cd04 SCHRAUSSER CODE, ASCII-combination (1 or 2 characters, e.g. -, #+, ..., minimum file size) 
               (5):HEX Hexadecimal value - No AIIB recoding
               (6):DEZ Decimal value
               (7):OKT Octal value – No AIIB recoding
               (8):ASCII characters directly - No AIIB recoding
[form] ....... (1):1-column ASCII output 
               (2):1-line ASCII output in block format without breaks
\end{verbatim}

Example \emph{{[}Bsp.{]}}:

\begin{verbatim}
BIIA bild.jpg bild.bii 4 2
\end{verbatim}

\hypertarget{aiib}{%
\subsubsection{5.10. AIIB}\label{aiib}}

\begin{center}\rule{0.5\linewidth}{0.5pt}\end{center}

Recodes a \emph{SCHRAUSSER CODE} or decimal encoded file (encoding via
\textbf{\emph{BIIA.EXE}}) \emph{{[}Recodiert eine SCHRAUSSER CODE oder
dezimal encodierte Datei (Encodierung via BIIA.EXE){]}}.

\begin{itemize}
\item
  Importing a \emph{SCHRAUSSER CODE} file (\textbf{\emph{*.bii}}).
\item
  Output of the recoded original file.
\item
  \emph{{[}Übernahme einer SCHRAUSSER CODE Datei (*.bii).{]}}
\item
  \emph{{[}Ausgabe der recodierten Ursprungsdatei.{]}}
\end{itemize}

Usage \emph{{[}Handhabung{]}}:

\begin{verbatim}
AIIB [input]
[input] ...... SCHRAUSSER CODE file
\end{verbatim}

Example \emph{{[}Bsp.{]}}:

\begin{verbatim}
AIIB bild.bii
\end{verbatim}

\hypertarget{stgersetz}{%
\subsubsection{5.11. STGERSETZ}\label{stgersetz}}

\begin{center}\rule{0.5\linewidth}{0.5pt}\end{center}

Modification of one or more \emph{substrings} (single characters or
\emph{strings}) in a contiguous \emph{string} (without spaces), read as
a line of an \emph{ASCII} file \emph{{[}Änderung ein oder mehrerer
Substrings (Einzelzeichen oder Zeichenketten) in einem zusammenhängenden
String (ohne Leerzeichen), eingelesen als Zeile einer ASCII Datei{]}}.

\begin{itemize}
\tightlist
\item
  Importing a line from an \emph{ASCII} file.
\item
  \emph{{[}Übernahme einer Zeile aus einem ASCII-File.{]}}
\end{itemize}

Example \emph{{[}Bsp.{]}}:

\begin{verbatim}
12345678
abcdefghijk
ABCDEFG
1234567890
\end{verbatim}

\begin{itemize}
\tightlist
\item
  Line-by-line processing of the \emph{strings} and output to an
  \emph{ASCII} file.
\item
  \emph{{[}zeilenweise Bearbeitung der Strings und Ausgabe in ein
  ASCII-File.{]}}
\end{itemize}

The processing steps are logged in \textbf{\emph{stgersetz\_log.txt}}
\emph{{[}Die Bearbeitungsschritte werden in stgersetz\_log.txt
protokolliert{]}}.

Usage \emph{{[}Handhabung{]}}:

\begin{verbatim}
STGERSETZ [input][output][typ][alt][neu]
[input] ...... Input file (strings line by line)
[output] ..... Output file (strings line by line)
[typ] ........ (0):deletes
               (1):modifies
               (2):insert a space
[alt] ........ String to be modified
[neu] ........ Change to string
\end{verbatim}

The type of argument {[}neu{]} has no effect on {[}typ{]} (0)
\emph{delete} and (2) \emph{insert a space}, but must be specified as
character (e.g.~`-'). A space in the input file splits the \emph{string}
and results in a new line in the output file. Some special characters
are ignored \emph{{[}Die Art des Arguments {[}neu{]} hat für {[}typ{]}
(0) löschen und (2) ein Leerzeichen einfügen keine Auswirkung, muss aber
als beliebiges Zeichen angegeben werden (iuF. `-'). Ein Leerzeichen in
der Eingabedatei trennt die Zeichenkette und führt zu einer neuen Zeile
in der Ausgabedatei. Manche Sonderzeichen bleiben unberücksichtigt{]}.}

Example \emph{{[}Bsp.{]}}:

\begin{verbatim}
STGERSETZ input.txt output.txt 0 8 - //deletes 8
\end{verbatim}

\begin{verbatim}
1234567
abcdefghijk
ABCDEFG
123456790
\end{verbatim}

\begin{verbatim}
STGERSETZ input.txt output.txt 1 def DEF //changes def to DEF
\end{verbatim}

\begin{verbatim}
12345678
abcDEFghijk
ABCDEFG
1234567890
\end{verbatim}

\begin{verbatim}
STGERSETZ input.txt output.txt 2 234 - //inserts a space for 234
\end{verbatim}

\begin{verbatim}
1 5678
abcdefghijk
ABCDEFG
1 567890
\end{verbatim}

\hypertarget{consoleapp_tools}{%
\subsection{6. ConsoleApp\_Tools}\label{consoleapp_tools}}

Console \emph{tool} applications, see Schrausser
(\href{https://doi.org/10.5281/zenodo.7655239}{2023e}).

\hypertarget{auto}{%
\subsubsection{6.1. AUTO}\label{auto}}

\begin{center}\rule{0.5\linewidth}{0.5pt}\end{center}

Creates a \emph{C}-source file and directory structure.

\begin{itemize}
\tightlist
\item
  Declaration of parmeters (\emph{name}, \emph{filestreams},
  \emph{loops}).
\item
  Creation of framework-file \textbf{\emph{NAME.c}} as specified.
\item
  Creation of \emph{ASCII} framework-file
  \textbf{\emph{NOTE\_NAME.txt}}.
\end{itemize}

Directory structure:

\begin{verbatim}
_NAME_\NAME.c
_NAME_\debug\NOTE_NAME.txt
             DOSCOM.exe
\end{verbatim}

Usage:

\begin{verbatim}
AUTO [name][input][output][nl][L]
[name] ....... C projekt mame 
[input] ...... ASCII input filename
[output] ..... ASCII output filename
[nl] ......... Number of loop-shells to implement
[L] .......... Language (0):English (1):German
\end{verbatim}

Example:

\begin{verbatim}
AUTO TEST 0 0 0 0
AUTO PRO1 input.txt output.txt 3 1
AUTO PRO2 0 output.txt 1 0
\end{verbatim}

\hypertarget{dls}{%
\subsubsection{6.2. DLS}\label{dls}}

\begin{center}\rule{0.5\linewidth}{0.5pt}\end{center}

Writes a file to the console \emph{{[}Schreibt eine Datei zur
Konsole{]}}.

Usage \emph{{[}Handhabung{]}}:

\begin{verbatim}
DLS [input] [ [i][m1][m2] ]
[input] ...... Input file 
optional:
[i] .......... (1):displays file name, n, k, m1, m2, i, j 
               (0):displays no additional information
[m1] ......... Start line of text output 
[m2] ......... End line of text output
\end{verbatim}

Example \emph{{[}Bsp.{]}}:

\begin{verbatim}
DLS NOTIZ_DLS.txt
DLS NOTIZ_DLS.txt 1
DLS NOTIZ_DLS.txt 1 1
DLS NOTIZ_DLS.txt 0 3 10
DLS NOTIZ_DLS.txt 1 10 10
\end{verbatim}

\hypertarget{fil}{%
\subsubsection{6.3. FIL}\label{fil}}

\begin{center}\rule{0.5\linewidth}{0.5pt}\end{center}

Writes a list of all specified files in the current directory to an
\emph{ASCII} file \emph{{[}Schreibt eine Liste aller spezifizierten
Dateien des aktuellen Verzeichnisses in eine ASCII Datei{]}}.

Usage \emph{{[}Handhabung{]}}:

\begin{verbatim}
FIL [output][*] 
[output] ..... Output file 
[*] .......... File specification
\end{verbatim}

Example \emph{{[}Bsp.{]}}:

\begin{verbatim}
FIL dir.txt *.*
\end{verbatim}

\hypertarget{kopie}{%
\subsubsection{6.4. KOPIE}\label{kopie}}

\begin{center}\rule{0.5\linewidth}{0.5pt}\end{center}

Copies a file \emph{{[}Kopiert eine Datei{]}}.

\begin{itemize}
\item
  Import of a file.
\item
  Output of a copy of a file (if applicable modified).
\item
  \emph{{[}Übernahme einer Datei.{]}}
\item
  \emph{{[}Ausgabe einer (ggf. modifizierten) Dateikopie.{]}}
\end{itemize}

Usage \emph{{[}Handhabung{]}}:

\begin{verbatim}
KOPIE [input][output] ([switch]) 
[input] ...... Input file 
[output] ..... Output file 
\end{verbatim}

Example \emph{{[}Bsp.{]}}:

\begin{verbatim}
KOPIE DTREN.BAT DTRENA.BAT
\end{verbatim}

\hypertarget{src}{%
\subsubsection{6.5. SRC}\label{src}}

\begin{center}\rule{0.5\linewidth}{0.5pt}\end{center}

Writes the command-line argument to an \emph{ASCII} file
\emph{{[}Schreibt das Kommandozeilenargument in eine ASCII Datei{]}}.

Usage \emph{{[}Handhabung{]}}:

\begin{verbatim}
SRC [output][typ][arg][sw]
[output] ..... Output file
[typ] ........ (0):attach (1):create 
[arg] ........ Argument
[sw] ......... (0):line break (1):space (2):no character
\end{verbatim}

Example \emph{{[}Bsp.{]}}:

\begin{verbatim}
SRC out.bat 1 lsn 1
SRC out.bat 0 nul.txt 0
\end{verbatim}

\hypertarget{t}{%
\subsubsection{6.6. T}\label{t}}

\begin{center}\rule{0.5\linewidth}{0.5pt}\end{center}

Inserts the commandline arguments into the \emph{ASCII} file
\textbf{\emph{T.txt}}. The arguments are separated by a space
\emph{{[}Fügt die Kommandozeilenargumente in die ASCII Datei T.txt ein.
Die Argumente werden in der Datei durch ein Leerzeichen getrennt{]}}.

Usage \emph{{[}Handhabung{]}}:

\begin{verbatim}
T [argi]…[argn]
[argi] ....... String
\end{verbatim}

Example \emph{{[}Bsp.{]}}:

\begin{verbatim}
T This is a very simple program
T for fast text input into a file.
T Das ist ein sehr einfaches Programm
T zur schnellen Texteingabe in eine Datei.
\end{verbatim}

\hypertarget{doscom}{%
\subsubsection{6.7. DOSCOM}\label{doscom}}

\begin{center}\rule{0.5\linewidth}{0.5pt}\end{center}

Starts a \emph{DOS} \emph{shell} in the current directory. Copies itself
to the directory \textbf{\emph{{[}\_DOSCOM{]}}}, which is linked to the
desktop. Renames to \emph{nil} after use in the current directory
\emph{{[}Startet eine DOS-Shell im aktuellen Verzeichnis. Kopiert sich
selbst in das mit dem Desktop verknüfte Verzeichis {[}\_DOSCOM{]}.
Umbenennung in nil nach Gebrauch im aktuellen Verzeichis{]}}.

Usage \emph{{[}Handhabung{]}}:

\emph{Drag} the desktop shortcut \textbf{\emph{{[}\_DOSCOM{]}}} to any
directory, run it there, and delete it after use \emph{{[}Aus der
Desktopverknüpfung {[}\_DOSCOM{]} in ein beliebiges Verzeichnis ziehen,
dort ausführen und nach Gebrauch löschen{]}}.

\hypertarget{finde}{%
\subsubsection{6.8. FINDE}\label{finde}}

\begin{center}\rule{0.5\linewidth}{0.5pt}\end{center}

Search and \emph{find} files \emph{{[}Dateien suchen und finden{]}}.

Usage \emph{{[}Handhabung{]}}:

\begin{verbatim}
FINDE [file]
[file] ....... File to be found
\end{verbatim}

\hypertarget{lsn}{%
\subsubsection{6.9. LSN}\label{lsn}}

\begin{center}\rule{0.5\linewidth}{0.5pt}\end{center}

Deletes a file \emph{{[}Löscht eine Datei{]}}.

\begin{itemize}
\tightlist
\item
  Input the name of a file to be deleted.
\item
  \emph{{[}Übernahme des Namens einer zu löschenden Datei.{]}}
\end{itemize}

Usage \emph{{[}Handhabung{]}}:

\begin{verbatim}
LSN [file] 
[file] ....... Filename 
\end{verbatim}

Example \emph{{[}Bsp.{]}}:

\begin{verbatim}
LSN oo
\end{verbatim}

\hypertarget{cmnd-cmndw}{%
\subsubsection{6.10. CMND, CMNDw}\label{cmnd-cmndw}}

\begin{center}\rule{0.5\linewidth}{0.5pt}\end{center}

Executes \emph{console} commands with arguments \emph{{[}Führt
Consolenbefehle mit Argumenten aus{]}}.

Usage \emph{{[}Handhabung{]}}:

\begin{verbatim}
CMND [cmd]
[cmd] ........ Command

CMNDw [cmd][argc]
[cmd] ........ Command
[argc] ....... n of Arguments
\end{verbatim}

\textbf{Table 3}. \emph{ConsoleApp} tools overview.

\begin{verbatim}
Category        chpt.   nr.     Function

String/ Number  5       1       2DEC
                5       2       DEC2
                5       3       BIN2DEC
                5       4       DEC2BIN
                5       5       ABCD
                5       6       AZUBE
                5       7       TEZUTE
                5       8       AB_L
                5       9       BIIA
                5       10      AIIB
                5       11      STGERSETZ

Tools           6       1       AUTO
                6       2       DLS
                6       3       FIL
                6       4       KOPIE
                6       5       SRC
                6       6       T
                6       7       DOSCOM
                6       8       FINDE
                6       9       LSN
                6       10      CMND, CMNDw           
\end{verbatim}

\hypertarget{references}{%
\subsection{References}\label{references}}

\end{document}
